% !TeX root = ../../Book1.tex
% 纲要标题：## 第16章　李群的局部结构与整体积分
% 纲要位置：几何学计划纲要.md，第 733 行。
% 正文状态：初稿；尚未正式审核。

% 纲要内容开始
% 
% > 本章目标与证明范围：在有限维实李群范围内建立指数、闭子群和同态积分。李代数到李群的局部与整体结论分开，证明单连通性消除的具体障碍。
% 
% 纲要内容结束

\chapter{李群的局部结构与整体积分}
\label{g1:ch:16}
\BookChapterNumber{g1:ch:16}

\begin{chapterabstract}
	李群把光滑变化与群乘法放在同一对象上，单位元切空间则记录这些变化的一阶方向。左平移将切向量延伸为向量场，其括号给出李代数。我们先构造指数映射，证明局部可逆性，并以 BCH 的有限阶展开解释非交换乘法怎样出现在切空间中。随后用 Frobenius 定理把李子代数积分成浸入子群，再证明闭子群必为嵌入子群。最后把李代数同态先积分到单位元附近，再沿路径延拓；连续同伦的网格论证说明单连通性恰好消除了绕行带来的歧义。万有覆盖群和离散中心商进一步说明，相同李代数可以对应不同的整体李群。
\end{chapterabstract}

\begin{learninggoals}
	\begin{itemize}
		\item 从矩阵例子与一般定义认识李群，构造单位元切空间上的李括号，证明伴随作用和同态微分的相容性。
		\item 证明左不变场完备，建立指数的局部逆、自然性及微分公式，推导 BCH 三阶展开和乘积极限。
		\item 利用左不变积分分布构造连通浸入子群，证明闭子群定理，并用无理斜率子群辨别叶拓扑与子空间拓扑。
		\item 从局部同态构造路径延拓，证明单连通源上的整体积分，构造万有覆盖群及离散中心商，并判断同态何时能下降到商群。
	\end{itemize}
\end{learninggoals}

旋转既可以连续改变角度，也可以前后复合；可逆线性变换也同时具有光滑参数和群运算。前面我们分别研究了变换群、流形及向量场，现在可以把这些结构放在一起。群乘法使每一点都能与单位元比较，因此看似遍布整个群的微分信息，有希望由一个切空间组织起来。

\ProseParagraphBreak

不过，一个切空间究竟能记住多少整体信息？实数加法群与圆周群的单位切空间都是一维、括号都是零，但沿同一个局部方向运动，前者不会绕回起点，后者会。反过来，环面上的一条无理斜率运动虽由一个实参数描述，其像却可以稠密。本章一边建立局部与整体之间的联系，一边保留周期、闭性和拓扑所造成的差别。

\ProseParagraphBreak

本章的李群均为有限维实光滑李群，流形无边界、Hausdorff、第二可数；复矩阵群在这里也先按实流形处理。读者需要\chapref{g1:ch:11}的逆函数、常秩及嵌入概念，\chapref{g1:ch:13}的常微分方程、李括号与 Frobenius 定理，以及\chapref{g1:ch:05}的路径同伦和万有覆盖。指数微分的计算还会回顾\secref{g1:sec:1401}的一形式外微分。每次调用时将指出实际条件，尤其区分浸入与嵌入、连通与单连通。

\ProseParagraphBreak

本章在正文证明指数、闭子群与同态积分。这里整体同态的存在定理以已给定的李群 \(G,H\) 为对象，并要求源群 \(G\) 单连通；任意抽象有限维李代数能否先找到一个李群来实现，是另一项存在问题，不包含在本章的证明范围内。我们也不预设抽象李群都可以忠实地表示为矩阵群。

\ProseParagraphBreak

延伸阅读可查 Lee 第二版第19章的 Lie Subalgebras and Lie Subgroups（第505--506页），以及第20章的 One-Parameter Subgroups and the Exponential Map（第516--521页）、The Closed Subgroup Theorem（第522--524页）和 The Lie Correspondence（第530--532页）\cite[第505--506、516--524、530--532页]{Lee2012Smooth}。矩阵计算可对照 Hall 第二版 §§2.1--2.4、3.3--3.5；指数微分、BCH、同态及覆盖的阅读位置为 §§5.3--5.9\cite[\S\S2.1--2.4, 3.3--3.5, 5.3--5.9]{Hall2015}。不同文献可能从矩阵群或抽象流形出发，对照时应核清对象范围及左、右不变场的约定。

% !TeX root = ../../Book1.tex
% 纲要标题：### 16.1 Lie 群与 Lie 代数
% 纲要位置：几何学计划纲要.md，第 737 行。
% 正文状态：初稿；尚未正式审核。

% 纲要内容开始
% 
% * 矩阵群的光滑结构
% * 左不变向量场和单位元切空间
% * 伴随作用、伴随表示与括号
% * Lie 群同态的微分
% 
% 纲要内容结束

\section{Lie 群与 Lie 代数}
\label{g1:sec:1601}

旋转、平移和可逆线性变换既能复合，又能随参数光滑变化。群结构记录复合的规则，流形结构允许我们对这些变化求导。两者相容之后，单位元附近的变化便能通过平移传播到整个群；原本分散在各点的切向量，可以统一放回单位元的切空间研究。不过，一个向量空间本身还没有记录变换复合的次序。我们将看到，左不变向量场的 Lie 括号恰好给这个切空间补上一种运算，使它保留局部复合的部分信息。

\ProseParagraphBreak

本节中的 Lie 群均为有限维实 Lie 群，底流形 Hausdorff、第二可数且无边界，不要求连通。即使元素写成复矩阵，其流形维数也先按实数计算。矩阵群提供可计算的模型，后面的定义和证明则适用于一般 Lie 群，不要求预先把群嵌入某个矩阵空间。

\subsection{矩阵群的光滑结构}
\label{g1:subsec:1601:01}

要对群运算求导，先要规定乘法与求逆怎样和流形结构相容。矩阵群提供了直接的检验方法：先由矩阵方程建立流形结构，再把乘法与求逆限制到这个流形上。

\begin{definition}[Lie 群]\label{g1:def:lie-group}
	一个群 \(G\) 若同时是有限维、Hausdorff、第二可数的无边界实光滑流形，并且乘法与求逆映射
	\[
		m:G\times G\longrightarrow G,\qquad m(g,h)=gh,
		\qquad
		\iota:G\longrightarrow G,\qquad \iota(g)=g^{-1}
	\]
	均光滑，就称为\term[liqun]{Lie 群}[Lie group]。群的单位元记为 \(e\)。两个 Lie 群之间的光滑群同态称为\term[liquntongtai]{Lie 群同态}[Lie group homomorphism]；若它还是微分同胚，就称为 Lie 群同构。
\end{definition}

例如 \((\mathbb R^n,+)\) 是 Lie 群，单位元为 \(0\)。有限个 Lie 群的直积取乘积流形结构后仍是 Lie 群，因为乘法和求逆逐分量进行。可数离散群取零维流形结构也是 Lie 群：所有点都是坐标邻域，所有映射都光滑，可数性则保证第二可数。这里允许零维和不连通的例子，是因为群的离散信息并不会自动消失，只是稍后的切空间无法记录全部这类信息。

\ProseParagraphBreak

对 \(g\in G\)，记左平移和右平移为
\[
	L_g(h)=gh,\qquad R_g(h)=hg.
\]
它们都是微分同胚，逆分别为 \(L_{g^{-1}},R_{g^{-1}}\)：光滑性来自固定乘法的一个变量，逆的公式来自群律。特别地，各点切空间都与 \(\mathrm{T}_eG\) 线性同构。乘法在单位元处的微分还有一个直接的表达式：由 \(m(g,e)=g\)、\(m(e,h)=h\) 及微分的线性，
\begin{equation}\label{g1:eq:lie-group-unit-differentials}
	\mathrm{d}m_{(e,e)}(X,Y)=X+Y,\qquad \mathrm{d}\iota_e(X)=-X.
\end{equation}
第二式是对恒等式 \(m(g,\iota(g))=e\) 求导的结果，所用链式法则已在流形微分的定义中建立。

\begin{proposition}[经典矩阵群]\label{g1:prop:matrix-classical-lie-groups}
	设 \(n\geq1\)。下表中的矩阵集合具有其所在实矩阵空间的自然开子流形或嵌入子流形结构，配上矩阵乘法都是实 Lie 群。这里 \(A^*=\overline A^{\mathsf T}\)，复矩阵空间按实向量空间理解；表中的切空间均在单位矩阵 \(I\) 处。
	\[
		\begin{array}{c|c|c}
			G&\dim_{\mathbb R}G&\mathrm{T}_IG\\ \hline
			\operatorname{GL}(n,\mathbb R)&n^2&\operatorname{Mat}_{n\times n}(\mathbb R)\\
			\operatorname{GL}(n,\mathbb C)&2n^2&\operatorname{Mat}_{n\times n}(\mathbb C)\\
			\operatorname{SL}(n,\mathbb R)&n^2-1&\{X:\operatorname{tr}X=0\}\\
			\operatorname{SL}(n,\mathbb C)&2n^2-2&\{X:\operatorname{tr}X=0\}\\
			\operatorname{O}(n),\ \operatorname{SO}(n)&n(n-1)/2&\{X:X^{\mathsf T}+X=0\}\\
			\operatorname{U}(n)&n^2&\{X:X^*+X=0\}\\
			\operatorname{SU}(n)&n^2-1&\{X:X^*+X=0,\ \operatorname{tr}X=0\}.
		\end{array}
	\]
	其中 \(\operatorname{GL}\) 由可逆矩阵组成，\(\operatorname{SL}\) 再要求行列式为 \(1\)；\(\operatorname{O}(n)=\{A\in\operatorname{GL}(n,\mathbb R):A^{\mathsf T}A=I\}\)，\(\operatorname{SO}(n)=\operatorname{O}(n)\cap\operatorname{SL}(n,\mathbb R)\)；\(\operatorname{U}(n)=\{A\in\operatorname{GL}(n,\mathbb C):A^*A=I\}\)，\(\operatorname{SU}(n)=\operatorname{U}(n)\cap\operatorname{SL}(n,\mathbb C)\)。表中每行的矩阵域与该行的群一致。
\end{proposition}
\begin{proof}
	行列式是矩阵分量的多项式，故非零行列式的集合 \(\operatorname{GL}(n,\mathbb R)\)、\(\operatorname{GL}(n,\mathbb C)\) 是开集，实维数分别为 \(n^2,2n^2\)。矩阵乘法的分量是多项式，求逆的公式为 \(A^{-1}=\operatorname{adj}(A)/\det A\)，因此两种运算都光滑。复矩阵情形也如此，因为复数的加法、乘法、共轭以及非零复数的求逆都是实光滑映射。

	先计算行列式的微分。按列的多线性展开 \(\det(I+tB)\)，常数项为 \(1\)，一次项只能从一列取 \(tB\)、其余列取单位矩阵的列，故一次项为 \(t\sum_jB_{jj}=t\operatorname{tr}B\)。于是对可逆 \(A\)，
	\begin{equation}\label{g1:eq:determinant-differential}
		\mathrm{d}(\det)_A(H)=\det A\,\operatorname{tr}(A^{-1}H).
	\end{equation}
	在实矩阵的 \(\det A=1\) 处，取 \(H=(c/n)A\) 就使微分等于任意 \(c\in\mathbb R\)；在复矩阵情形，取 \(H=(z/n)A\) 则得到任意 \(z\in\mathbb C\)。因此 \(1\) 分别是到 \(\mathbb R\) 和到 \(\mathbb C\simeq\mathbb R^2\) 的行列式映射的正则值。由\cref{g1:thm:regular-level-set}，两种 \(\operatorname{SL}\) 都是嵌入子流形，其余维分别为 \(1,2\)，单位元切空间均由 \(\operatorname{tr}X=0\) 给出。

	对正交群，令 \(\operatorname{Sym}_n(\mathbb R)\) 为实对称矩阵空间，其维数为 \(n(n+1)/2\)，考虑
	\[
		Q:\operatorname{GL}(n,\mathbb R)\longrightarrow\operatorname{Sym}_n(\mathbb R),\qquad Q(A)=A^{\mathsf T}A.
	\]
	这里把靶空间取成对称矩阵空间很重要：若仍取全体矩阵，微分便不可能满射。在 \(Q(A)=I\) 处，
	\[
		\mathrm{d}Q_A(H)=A^{\mathsf T}H+H^{\mathsf T}A.
	\]
	给定对称矩阵 \(S\)，取 \(H=AS/2\)，就有 \(\mathrm{d}Q_A(H)=S\)。故 \(I\) 是正则值，\(\operatorname{O}(n)\) 是维数 \(n^2-n(n+1)/2=n(n-1)/2\) 的嵌入子流形，其单位元切空间为反对称矩阵。正交矩阵满足 \((\det A)^2=1\)，所以 \(\operatorname{SO}(n)\) 正是 \(\operatorname{O}(n)\) 中 \(\det A>0\) 的开子集，具有相同的局部维数和单位元切空间。

	复矩阵情形令 \(\operatorname{Herm}_n(\mathbb C)=\{S:S^*=S\}\)。它的对角项有 \(n\) 个实自由度，上三角有 \(n(n-1)/2\) 个复自由度，实维数为 \(n^2\)。映射 \(Q(A)=A^*A\) 的微分为 \(A^*H+H^*A\)。在 \(A^*A=I\) 处，对任意 Hermitian 矩阵 \(S\)，\(H=AS/2\) 仍给出 \(\mathrm{d}Q_A(H)=S\)。故 \(\operatorname{U}(n)\) 是实维数 \(2n^2-n^2=n^2\) 的嵌入子流形，\(\mathrm{T}_I\operatorname{U}(n)\) 由 \(X^*+X=0\) 给出。

	最后讨论 \(\operatorname{SU}(n)\)。圆周 \(S^1=\{z\in\mathbb C:|z|^2=1\}\) 是一维嵌入子流形，因为 \(z\mapsto|z|^2\) 在圆周上微分非零，其切空间为 \(\mathrm{T}_zS^1=iz\mathbb R\)。酉矩阵的行列式属于 \(S^1\)，由\cref{g1:prop:embedded-initial-property}，行列式限制给出光滑映射 \(\det:\operatorname{U}(n)\to S^1\)。上述 \(Q\) 的微分还说明
	\[
		\mathrm{T}_A\operatorname{U}(n)=\{AB:B^*+B=0\}.
	\]
	由\eqref{g1:eq:determinant-differential}，\(\mathrm{d}(\det)_A(AB)=\det A\,\operatorname{tr}B\)。取 \(B=(it/n)I\)，可得到 \(i\det A\,t\)，遍历 \(\mathrm{T}_{\det A}S^1\)。于是 \(1\) 是这次行列式映射的正则值，\(\operatorname{SU}(n)\) 的实余维为 \(1\)，单位元切空间再加上 \(\operatorname{tr}X=0\) 这一条件。反 Hermitian 矩阵的迹本来就是纯虚数，因此这里的迹零只有一个实条件，与 \(\operatorname{SL}(n,\mathbb C)\) 中的两个实条件并不矛盾。

	各集合由定义直接对乘法和求逆封闭。它们在上述构造中都具有嵌入子流形结构；环境矩阵群的乘法和求逆限制后仍光滑地取值于相应集合，理由再次是\cref{g1:prop:embedded-initial-property}。这些子流形的拓扑是子空间拓扑，因此 Hausdorff、第二可数，无边界性由正则值定理或开子流形构造保证。
\end{proof}

以上矩阵群的流形结构来自方程微分的秩。例如，正交群的证明把对称矩阵作为靶空间，使单位矩阵成为正则值。稍后的一般闭子群定理将从子群的闭性统一构造局部坐标，给出另一种建立流形结构的方法。

\subsection{左不变向量场和单位元切空间}
\label{g1:subsec:1601:02}

现在设 \(G\) 是任意有限维实 Lie 群，记 \(\mathfrak g=\mathrm{T}_eG\)。一个方向 \(X\in\mathfrak g\) 可以被左平移送到每个点：在 \(g\) 处取 \(\mathrm{d}(L_g)_eX\)。这个向量场精确表达了“各点都朝同一个群方向前进”。它一般不是某个坐标系中的常向量场；保持不变的是左平移给出的比较方式。

\begin{proposition}[左不变场]\label{g1:prop:left-invariant-fields}
	设 \(G\) 为有限维实 Lie 群。若 \(Z\in\mathfrak X(G)\) 对任意 \(g,h\in G\) 都满足
	\[
		Z_{gh}=\mathrm{d}(L_g)_hZ_h.
	\]
	则称它为\term[zuobubianxiangliangchang]{左不变向量场}[left-invariant vector field]。在单位元处求值，得到左不变向量场空间到 \(\mathfrak g=\mathrm{T}_eG\) 的实线性同构，其逆为
	\begin{equation}\label{g1:eq:left-invariant-extension}
		X\longmapsto X^L,\qquad X^L_g=\mathrm{d}(L_g)_eX.
	\end{equation}
	两个左不变场的 Lie 括号仍左不变。\(\mathfrak g\) 的任意一组基经上述延拓，得到 \(\mathrm{T}G\) 的一个整体光滑标架。
\end{proposition}
\begin{proof}
	乘法 \(m(g,h)=gh\) 光滑，对第二个变量在 \(h=e\) 处求微分，得到随 \(g\) 光滑变化的线性映射 \(\mathrm{d}(L_g)_e\)。因而 \(X^L\) 确为光滑向量场。由 \(L_{gh}=L_g\circ L_h\) 及链式法则，
	\[
		X^L_{gh}=\mathrm{d}(L_g)_h\,\mathrm{d}(L_h)_eX=\mathrm{d}(L_g)_hX^L_h,
	\]
	所以它左不变。反过来，若 \(Z\) 左不变，在定义中令 \(h=e\)，立即得到 \(Z_g=\mathrm{d}(L_g)_eZ_e\)，故 \(Z\) 完全由 \(Z_e\) 确定。两种操作互逆且都线性。

	左不变等价于 \((L_g)_*Z=Z\)。由\cref{g1:prop:lie-derivative-natural}中微分同胚保持 Lie 括号的结论，若 \(Z,W\) 左不变，则
	\[
		(L_g)_*[Z,W]=[(L_g)_*Z,(L_g)_*W]=[Z,W].
	\]
	故其括号也左不变。若 \(X_1,\ldots,X_d\) 为 \(\mathfrak g\) 的基，\(\mathrm{d}(L_g)_e\) 是线性同构，故 \(X_1^L(g),\ldots,X_d^L(g)\) 在每点构成基，给出整体标架。具体的切丛平凡化为 \((g,X)\mapsto\mathrm{d}(L_g)_eX\)，其逆在 \(\mathrm{T}_gG\) 上为 \(v\mapsto(g,\mathrm{d}(L_{g^{-1}})_gv)\)；乘法、求逆及其微分的光滑性说明正逆映射都光滑。
\end{proof}

\begin{definition}[Lie 代数]\label{g1:def:lie-algebra-group}
	实向量空间 \(V\) 上一个实双线性运算 \([\ ,\ ]:V\times V\to V\)，若满足反对称性 \([X,Y]=-[Y,X]\) 及 Jacobi 恒等式
	\[
		[X,[Y,Z]]+[Y,[Z,X]]+[Z,[X,Y]]=0,
	\]
	就使 \(V\) 成为一个\term[lidaishu]{Lie 代数}[Lie algebra]。对有限维实 Lie 群 \(G\)，在 \(\mathfrak g=\mathrm{T}_eG\) 上定义
	\[
		[X,Y]\coloneqq[X^L,Y^L]_e.
	\]
	所得 Lie 代数称为 \(G\) 的 Lie 代数，记为 \(\operatorname{Lie}(G)=\mathfrak g\)。\symidx{\(\mathfrak g=\mathrm{T}_eG\)：Lie 群的 Lie 代数}由\cref{g1:prop:left-invariant-fields}，其左不变延拓满足 \(([X,Y])^L=[X^L,Y^L]\)；实双线性、反对称性和 Jacobi 恒等式随即由\cref{g1:prop:lie-bracket-properties}的向量场版本传下来。
\end{definition}

一般向量场的 Lie 括号在一点的值还依赖附近的变化，不能只由该点的两个切向量计算。这里能够定义 \([X,Y]\in \mathrm{T}_eG\)，并非删掉了这种依赖，而是由左不变性为每个单位元切向量选定了唯一的整体延拓。因此群乘法参与了括号的定义；只给同一个流形选一个点，通常没有这样的运算。

\ProseParagraphBreak

以后还会用到两个代数概念。Lie 代数之间保持括号的线性映射称为\term[lidaishutongtai]{Lie 代数同态}[Lie algebra homomorphism]；线性子空间 \(\mathfrak h\subset\mathfrak g\) 若满足 \([\mathfrak h,\mathfrak h]\subset\mathfrak h\)，称为\term[lizidaishu]{Lie 子代数}[Lie subalgebra]。若更强地满足 \([\mathfrak g,\mathfrak h]\subset\mathfrak h\)，则称 \(\mathfrak h\) 为理想。前者保证子空间内部可以作括号，后者还控制它与整个 Lie 代数的相互作用。

\begin{proposition}[矩阵括号]\label{g1:prop:matrix-lie-bracket}
	设 \(n\ge1\)，\(\mathbb F=\mathbb R\) 或 \(\mathbb C\)，\(G\) 为 \(\operatorname{GL}(n,\mathbb F)\)、\(\operatorname{SL}(n,\mathbb F)\)、\(\operatorname{O}(n)\)、\(\operatorname{SO}(n)\)、\(\operatorname{U}(n)\)、\(\operatorname{SU}(n)\) 之一。这些矩阵群的流形结构见\cref{g1:prop:matrix-classical-lie-groups}。在各群相应的矩阵域中，将 \(\mathfrak g=\mathrm{T}_IG\) 视为矩阵空间的实线性子空间。则 Lie 代数的括号为
	\[
		[X,Y]=XY-YX.
	\]
\end{proposition}
\begin{proof}
	先在 \(\operatorname{GL}(n,\mathbb R)\) 或 \(\operatorname{GL}(n,\mathbb C)\) 中计算。左平移 \(L_A(B)=AB\) 的微分为 \(H\mapsto AH\)，故 \(X^L_A=AX\)、\(Y^L_A=AY\)。把它们看成开矩阵空间上的向量值函数，\(\mathrm{d}Y^L_A(H)=HY\)、\(\mathrm{d}X^L_A(H)=HX\)。依\eqref{g1:eq:lie-bracket-coordinates}的符号约定，
	\[
		[X^L,Y^L]_A
		=\mathrm{d}Y^L_A(AX)-\mathrm{d}X^L_A(AY)
		=A(XY-YX).
	\]
	在 \(A=I\) 处即得所需公式。复矩阵情况下以上乘法仍是实线性微分，计算不变。

	其余矩阵群 \(H\) 的左平移仍由矩阵乘法给出。若 \(X\in \mathrm{T}_IH\)，\(A\mapsto AX\) 在 \(H\) 上取值于 \(\mathrm{T}H\)，且通过嵌入 \(H\hookrightarrow\operatorname{GL}\) 与环境的同名左不变场相关。再次用\cref{g1:prop:lie-derivative-natural}，它们的括号也与环境括号相关，所以单位元处的公式仍为 \(XY-YX\)。这也从几何上证明了表中各切空间对矩阵交换子封闭。
\end{proof}

这些 Lie 代数通常记作 \(\mathfrak{gl}(n,\mathbb R)\)、\(\mathfrak{sl}(n,\mathbb R)\)、\(\mathfrak{so}(n)\)、\(\mathfrak u(n)\)、\(\mathfrak{su}(n)\) 等；\(\operatorname{O}(n)\) 和 \(\operatorname{SO}(n)\) 具有同一个 \(\mathfrak{so}(n)\)。\(\mathfrak{gl}(n,\mathbb C)\)、\(\mathfrak{sl}(n,\mathbb C)\) 此处仍可视为实 Lie 代数，只是它们还额外具有与括号相容的复向量空间结构。

\ProseParagraphBreak

Lie 代数会丢失一些整体信息。例如 \((\mathbb R,+)\) 和 \(S^1=\operatorname{U}(1)\) 的 Lie 代数都是一维实向量空间，其括号因反对称性而恒为零；但两群不可能同构，因为圆周紧而实直线不紧。任意离散 Lie 群的 Lie 代数甚至都是零空间。因此本章后面既要研究局部积分，也要另行考虑覆盖、单连通性和离散子群。

\subsection{伴随作用、伴随表示与括号}
\label{g1:subsec:1601:03}

固定 \(g\in G\)，共轭 \(C_g(h)=ghg^{-1}\) 把单位元固定，因而其微分成为 \(\mathfrak g\) 的一个线性变换。左平移把单位元处的方向搬到 \(g\)；共轭则将改变后的方向仍放在单位元处比较，由此得到伴随表示。

\begin{proposition}[伴随表示]\label{g1:prop:adjoint-representation}
	设 \(G\) 为有限维实 Lie 群，\(\mathfrak g=\mathrm{T}_eG\)。对 \(g\in G\)，令共轭映射 \(C_g:G\to G\) 为 \(C_g(h)=ghg^{-1}\)，并定义
	\[
		\operatorname{Ad}_g\coloneqq\mathrm{d}(C_g)_e:\mathfrak g\longrightarrow\mathfrak g.
	\]
	则 \(\operatorname{Ad}:G\to\operatorname{GL}(\mathfrak g)\) 是光滑群同态，且每个 \(\operatorname{Ad}_g\) 保持 Lie 括号。它称为 \(G\) 的\term[bansuibiaoshi]{伴随表示}[adjoint representation]。这里一般称一个到有限维向量空间可逆线性变换群的光滑群同态为该 Lie 群的表示。
\end{proposition}
\begin{proof}
	共轭 \(C_g\) 是微分同胚，逆为 \(C_{g^{-1}}\)，因此 \(\operatorname{Ad}_g\) 可逆。由 \(C_{gh}=C_g\circ C_h\) 及它们都固定 \(e\)，链式法则给出
	\[
		\operatorname{Ad}_{gh}=\operatorname{Ad}_g\operatorname{Ad}_h,\qquad
		\operatorname{Ad}_e=\operatorname{id}_{\mathfrak g}.
	\]
	映射 \((g,h)\mapsto ghg^{-1}\) 光滑，对第二个变量在 \(h=e\) 处求微分，得到的矩阵分量随 \(g\) 光滑。这说明 \(\operatorname{Ad}\) 是到 \(\operatorname{End}(\mathfrak g)\) 的光滑映射，而其值落在开子集 \(\operatorname{GL}(\mathfrak g)\)，所以也是到后者的光滑映射。

	恒等式 \(C_g\circ L_h=L_{C_g(h)}\circ C_g\) 在单位元处求导，给出
	\[
		(C_g)_*X^L=(\operatorname{Ad}_gX)^L.
	\]
	微分同胚保持向量场的 Lie 括号，故
	\[
		(C_g)_*[X^L,Y^L]
		=[(\operatorname{Ad}_gX)^L,(\operatorname{Ad}_gY)^L].
	\]
	两边在 \(e\) 处求值，即得 \(\operatorname{Ad}_g[X,Y]=[\operatorname{Ad}_gX,\operatorname{Ad}_gY]\)。
\end{proof}

再对 \(\operatorname{Ad}\) 在单位元处求微分。利用 \(\mathrm{T}_I\operatorname{GL}(\mathfrak g)=\operatorname{End}(\mathfrak g)\)，记
\[
	\operatorname{ad}\coloneqq\mathrm{d}(\operatorname{Ad})_e:
	\mathfrak g\longrightarrow\operatorname{End}(\mathfrak g),\qquad
	X\longmapsto\operatorname{ad}_X.
\]
\symidx{\(\operatorname{Ad}_g,\operatorname{ad}_X\)：伴随表示及其微分}大写 \(\operatorname{Ad}\) 的输入是群元素，小写 \(\operatorname{ad}\) 的输入是单位元切向量。下面的公式把这次求导与此前定义的 Lie 括号联系起来。

\begin{proposition}[伴随微分]\label{g1:prop:adjoint-differential-bracket}
	设 \(G\) 为有限维实 Lie 群，\(\mathfrak g=\mathrm{T}_eG\)，\(\operatorname{Ad}:G\to\operatorname{GL}(\mathfrak g)\) 为伴随表示。在 \(\mathrm{T}_I\operatorname{GL}(\mathfrak g)=\operatorname{End}(\mathfrak g)\) 的识别下，令 \(\operatorname{ad}=(\mathrm{d}\operatorname{Ad})_e\)。则对任意 \(X,Y\in\mathfrak g\)，
	\begin{equation}\label{g1:eq:ad-bracket}
		\operatorname{ad}_XY=[X,Y].
	\end{equation}
	此外 \(\operatorname{ad}_X\) 是括号的导子，而且
	\[
		[\operatorname{ad}_X,\operatorname{ad}_Y]=\operatorname{ad}_{[X,Y]},
	\]
	左边是线性自同态的交换子。
\end{proposition}
\begin{proof}
	由\cref{g1:thm:maximal-integral-curve}，左不变场 \(X^L\) 在单位元有局部积分曲线 \(\gamma(t)\)，满足 \(\gamma(0)=e\)、\(\gamma'(0)=X\)。只需在 \(0\) 附近取它，不需要知道它是否能延拓到所有时间。对固定的 \(g\)，左不变性给出
	\[
		\frac{\mathrm{d}}{\mathrm{d}t}\bigl(g\gamma(t)\bigr)
		=\mathrm{d}(L_g)_{\gamma(t)}X^L_{\gamma(t)}
		=X^L_{g\gamma(t)}.
	\]
	它在 \(t=0\) 经过 \(g\)，所以由积分曲线唯一性，\(X^L\) 的局部流为
	\begin{equation}\label{g1:eq:left-invariant-local-flow}
		\Phi_t(g)=g\gamma(t)=R_{\gamma(t)}(g)
	\end{equation}
	在两边定义的局部时间范围内成立。因而流的逆局部映射是 \(R_{\gamma(t)^{-1}}\)。注意，左不变场的流在这里以右乘写出；这是 \(g\) 固定而单位元曲线被左平移的结果。

	按照\cref{g1:prop:flow-pullback-derivative}的向量场拉回定义，在单位元有
	\begin{align*}
		(\Phi_t^*Y^L)_e
		&=\mathrm{d}(R_{\gamma(t)^{-1}})_{\gamma(t)}Y^L_{\gamma(t)}\\
		&=\mathrm{d}(R_{\gamma(t)^{-1}})_{\gamma(t)}\,\mathrm{d}(L_{\gamma(t)})_eY\\
		&=\mathrm{d}(C_{\gamma(t)})_eY
		=\operatorname{Ad}_{\gamma(t)}Y.
	\end{align*}
	对 \(t\) 在 \(0\) 处求导。左边由流拉回的导数公式等于 \([X^L,Y^L]_e=[X,Y]\)，右边由 \(\gamma'(0)=X\) 等于 \(\operatorname{ad}_XY\)。这证明\eqref{g1:eq:ad-bracket}，也固定了它与第13章括号约定的符号。

	最后，Jacobi 恒等式改写为
	\[
		[X,[Y,Z]]=[[X,Y],Z]+[Y,[X,Z]],
	\]
	恰好说明 \(\operatorname{ad}_X[Y,Z]=[\operatorname{ad}_XY,Z]+[Y,\operatorname{ad}_XZ]\)，即导子性质。将第二项移到左边，又得到 \((\operatorname{ad}_X\operatorname{ad}_Y-\operatorname{ad}_Y\operatorname{ad}_X)Z=\operatorname{ad}_{[X,Y]}Z\)。对所有 \(Z\) 成立便是所需的自同态恒等式。
\end{proof}

对矩阵群，共轭的微分直接给出 \(\operatorname{Ad}_AY=AYA^{-1}\)。若 \(A(t)\) 是经过 \(I\)、速度为 \(X\) 的群内曲线，对 \(A(t)A(t)^{-1}=I\) 求导可得 \((A^{-1})'(0)=-X\)，从而 \(\frac{\mathrm{d}}{\mathrm{d}t}|_0 A(t)YA(t)^{-1}=XY-YX\)，与抽象证明一致。

\subsection{Lie 群同态的微分}
\label{g1:subsec:1601:04}

群同态保持乘法，因而应当保持由乘法定义的无穷小运算。左不变场把这个问题转为向量场的比较：同态的微分先把它们变成相关向量场，再利用第13章已经证明的括号自然性。

\begin{proposition}[同态的微分]\label{g1:prop:lie-homomorphism-differential}
	设 \(F:G\to H\) 是有限维实 Lie 群之间的 Lie 群同态，\(\mathfrak g=\mathrm{T}_{e_G}G\)、\(\mathfrak h=\mathrm{T}_{e_H}H\)。则
	\[
		\varphi\coloneqq\mathrm{d}F_{e_G}:\mathfrak g\longrightarrow\mathfrak h
	\]
	是 Lie 代数同态。具体地，对任意 \(g\in G\)、\(X,Y\in\mathfrak g\)，
	\begin{align*}
		\mathrm{d}F_g(X^L_g)&=(\varphi X)^L_{F(g)},\\
		\varphi[X,Y]&=[\varphi X,\varphi Y],\\
		\varphi\circ\operatorname{Ad}_g&=\operatorname{Ad}_{F(g)}\circ\varphi.
	\end{align*}
	同态复合的单位元微分为单位元微分的复合，恒等同态的微分为恒等映射。特别地，\(\operatorname{im}\varphi\) 是 \(\mathfrak h\) 的 Lie 子代数，\(\ker\varphi\) 是 \(\mathfrak g\) 的理想。
\end{proposition}
\begin{proof}
	群同态满足 \(F(e_G)=e_H\) 及 \(F\circ L_g=L_{F(g)}\circ F\)。后一式在 \(e_G\) 处求导，得到
	\[
		\mathrm{d}F_g\,\mathrm{d}(L_g)_{e_G}X
		=\mathrm{d}(L_{F(g)})_{e_H}\,\mathrm{d}F_{e_G}X,
	\]
	即第一式。因此 \(X^L\) 与 \((\varphi X)^L\) 是 \(F\)-相关场，\(Y^L\) 与 \((\varphi Y)^L\) 也如此。\cref{g1:prop:lie-derivative-natural}说明它们的括号仍相关。这一结论无需 \(F\) 可逆或满射：对 \(H\) 上任意局部光滑函数 \(f\)，只要把两次求导展开，就有
	\[
		[X^L,Y^L](f\circ F)
		=\bigl([(\varphi X)^L,(\varphi Y)^L]f\bigr)\circ F.
	\]
	在 \(e_G\) 处求值，恰得 \(\varphi[X,Y]=[\varphi X,\varphi Y]\)。同态还满足 \(F\circ C_g=C_{F(g)}\circ F\)，在单位元求导便给出伴随表示的交织公式。

	复合与恒等映射的断言是链式法则。若 \(U=\varphi X\)、\(V=\varphi Y\)，则 \([U,V]=\varphi[X,Y]\) 仍在像中；若 \(Y\in\ker\varphi\)，则对任意 \(X\in\mathfrak g\)，\(\varphi[X,Y]=[\varphi X,0]=0\)。这分别证明像是 Lie 子代数、核是理想。
\end{proof}

\begin{proposition}[同态的秩与核]\label{g1:prop:lie-homomorphism-constant-rank}
	设 \(F:G\to H\) 是有限维实 Lie 群之间的 Lie 群同态，令 \(r=\operatorname{rank}\mathrm{d}F_{e_G}\)。则 \(F\) 在所有点的秩都等于 \(r\)。核
	\[
		K=\ker F=F^{-1}(e_H)
	\]
	是 \(G\) 的闭嵌入子流形，维数为 \(\dim G-r\)。它以限制的群运算成为 Lie 群，且
	\[
		\mathrm{T}_{e_G}K=\ker\mathrm{d}F_{e_G};
	\]
	这里用包含映射的微分将左边视为 \(\mathrm{T}_{e_G}G\) 的子空间。
\end{proposition}
\begin{proof}
	前一命题的第一式可以写成
	\[
		\mathrm{d}F_g
		=\mathrm{d}(L_{F(g)})_{e_H}\circ\mathrm{d}F_{e_G}\circ\mathrm{d}(L_{g^{-1}})_g.
	\]
	两端的平移微分都是同构，故各点的秩都等于 \(r\)。固定 \(k\in K\)，由\cref{g1:cor:manifold-constant-rank}，在 \(k\) 和 \(e_H\) 附近可取分别以这两点为原点的坐标，使 \(F\) 写成
	\[
		(u^1,\ldots,u^m)\longmapsto(u^1,\ldots,u^r,0,\ldots,0),\qquad m=\dim G.
	\]
	必要时缩小源坐标邻域，使其像包含于选定的靶坐标邻域。于是核在这个源邻域中恰由 \(u^1=\cdots=u^r=0\) 给出，而 \((u^{r+1},\ldots,u^m)\) 是核的局部坐标。坐标转换由环境转换限制到这些切片得到，光滑且互逆。这些局部描述与子空间拓扑一致，给出 \(K\) 的嵌入子流形结构；它从 \(G\) 继承 Hausdorff 性和第二可数性。其切空间在每个 \(k\in K\) 都是 \(\ker\mathrm{d}F_k\)，因为坐标中两者均由前 \(r\) 个分量为零的向量组成。

	单点 \(\{e_H\}\) 在 Hausdorff 空间 \(H\) 中闭，\(F\) 连续，故 \(K\) 闭。群同态的核是子群，因此 \(G\) 的乘法和求逆限制后取值于 \(K\)。这些限制作为到 \(G\) 的映射光滑，由\cref{g1:prop:embedded-initial-property}，作为到 \(K\) 的映射也光滑。由此 \(K\) 满足 Lie 群的全部条件，单位元切空间的公式已在局部坐标计算中得到。
\end{proof}

恒秩定理把核写成坐标切片，给出它的嵌入子流形结构；如果 \(\mathrm{d}F_{e_G}\) 不满射，\(e_H\) 并非正则值，这时正是恒秩标准形发挥作用。像的整体拓扑还受远处源点的累积影响，\(F(G)\) 在 \(H\) 中是否闭或嵌入，需要下一节对子群作进一步分析。

\ProseParagraphBreak

例如实矩阵群的行列式同态 \(\det:\operatorname{GL}(n,\mathbb R)\to\mathbb R^\times\) 在单位元的微分是迹，核是 \(\operatorname{SL}(n,\mathbb R)\)，与此前的直接计算一致。若一般同态的单位元微分单射，它在各点都是浸入；但它仍可能把不同群元素送到同一点。下一节的指数映射将进一步把单位元切向量与群内曲线联系起来，之后才能系统讨论一个 Lie 代数同态能够恢复多少群同态的信息。

% 纲要条目排为编号小节标题，并进入目录与书签；说明段仍保留于纲要注释中。

% !TeX root = ../../Book1.tex
% 纲要标题：### 16.2 指数映射与局部乘法
% 纲要位置：几何学计划纲要.md，第 744 行。
% 正文状态：初稿；尚未正式审核。

% 纲要内容开始
% 
% * 单参数子群及完备性
% * 指数映射的微分和局部逆
% * 共轭与指数的相容
% * 局部乘法与 BCH 的三阶展开
% 
% 纲要内容结束

\section{指数映射与局部乘法}
\label{g1:sec:1602}

一个单位元处的切向量描述瞬时运动。左平移把同一个瞬时方向放到每一点，形成左不变向量场；沿它运动，就有望把切向量变成群元素。这条路线给出指数映射。它还回答另一个问题：把两个很小的运动相乘，结果与将两个切向量直接相加有多大差别？Lie 括号将成为最先出现的修正项。

\ProseParagraphBreak

本节的 Lie 群均为有限维实 Lie 群，流形无边界、Hausdorff、第二可数。记单位元为 \(e\)，Lie 代数为 \(\mathfrak g=\mathrm{T}_eG\)，\(X\in\mathfrak g\) 对应的左不变向量场记为 \(X^L_g=(\mathrm{d}L_g)_eX\)。这一对应及其括号相容性已在\secref{g1:sec:1601}建立。以下所有群乘积都有顺序，不能先把它们当作向量加法。

\subsection{单参数子群及完备性}
\label{g1:subsec:1602:01}

\begin{definition}\label{g1:def:one-parameter-subgroup}
	设 \(G\) 为有限维实 Lie 群。光滑同态 \(c:(\mathbb R,+)\to G\) 称为\term[dancanshuziqun]{单参数子群}[one-parameter subgroup]，即
	\[
		c(s+t)=c(s)c(t),\qquad c(0)=e.
	\]
	这里名称指带参数的同态，不要求 \(c\) 单射，也不预设其像是嵌入子流形。
\end{definition}

若 \(X=c'(0)\)，对 \(c(t+h)=c(t)c(h)\) 在 \(h=0\) 求导，便得 \(c'(t)=(\mathrm{d}L_{c(t)})_eX\)。所以单参数子群必是从单位元出发的左不变场积分曲线。反过来，局部初值定理只给出一小段积分曲线；要获得整个 \(\mathbb R\) 上的同态，还须证明它不会在有限时间终止。

\begin{theorem}[指数映射]\label{g1:thm:lie-exponential}
	设 \(G\) 为有限维实 Lie 群。对每个 \(X\in\mathfrak g\)，左不变场 \(X^L\) 完备，且存在唯一单参数子群 \(c_X\) 满足 \(c_X'(0)=X\)。定义
	\[
		\exp_G:\mathfrak g\longrightarrow G,\qquad \exp_GX=c_X(1).
	\]
	则 \(\exp_G\) 光滑，并且
	\begin{equation}\label{g1:eq:one-parameter-exponential}
		c_X(t)=\exp_G(tX),\qquad
		\exp_G((s+t)X)=\exp_G(sX)\exp_G(tX).
	\end{equation}
	\(X^L\) 的整体流为 \(\Phi_t(g)=g\exp_G(tX)\)。
\end{theorem}
\begin{proof}
	由\cref{g1:thm:maximal-integral-curve}先取过 \(e\) 的局部积分曲线 \(c\)，将其限制在 \((-2a,2a)\)，其中 \(a>0\)。左不变性说明，对于固定 \(s\)，\(t\mapsto c(s)c(t)\) 与 \(t\mapsto c(s+t)\) 都是 \(X^L\) 的积分曲线，并在 \(t=0\) 取相同值。因此只要相关参数留在这段局部区间中，唯一性给出局部乘法规律 \(c(s)c(t)=c(s+t)\)。

	现在把同一小段反复左平移。对每个整数 \(n\)，在区间 \((na-a,na+a)\) 上令
	\[
		\widehat c(t)=c(a)^n c(t-na).
	\]
	相邻区间交叠时，令 \(u=t-(n+1)a\in(-a,0)\)，则 \(c(a)c(u)=c(a+u)\)，所以两式相同。每段是光滑积分曲线，因而粘成整个实轴上的光滑积分曲线，且在零点取 \(e\)。任意 \(g\in G\) 出发的曲线 \(g\widehat c(t)\) 也在整个实轴有定义，所以 \(X^L\) 完备。对全局曲线再次使用初值唯一性，得到 \(\widehat c(s+t)=\widehat c(s)\widehat c(t)\)。前段对任意单参数子群的求导说明，另一个具有相同初速度的单参数子群必须是同一初值解，故唯一。记它为 \(c_X\)。

	曲线 \(s\mapsto c_X(ts)\) 满足场 \((tX)^L\) 的方程，并在零点取单位元，包括 \(t=0\) 和负 \(t\) 的情形。因此唯一性给出 \(c_{tX}(s)=c_X(ts)\)，令 \(s=1\) 便得 \(\exp_G(tX)=c_X(t)\)。对 \(X\) 的光滑性调用\cref{g1:thm:ode-smooth-dependence}：场 \((g,X)\mapsto(\mathrm{d}L_g)_eX\) 联合光滑，任意固定 \(X_0\) 的轨道已在紧时间区间 \([0,1]\) 上存在，所以附近 \(X\) 的解在共同区间上光滑依赖参数。于 \(t=1\) 求值，得到 \(\exp_G\) 在每个 \(X_0\) 附近光滑。最后，从 \(g\) 出发的积分曲线已经写成 \(g c_X(t)\)，即所列流公式。
\end{proof}

这个映射称为 Lie 群的\term[zhishuyingshe]{指数映射}[exponential map]。\symidx{\(\exp_G\)：Lie 群的指数映射}它的定义依赖群结构和左不变场，与第九章由度量和测地线定义的指数映射具有不同的输入数据。两者都把切向量送到某条指定曲线的终点，但不能未经说明互换。

\begin{example}\label{g1:exm:matrix-exponential-series}
	在 \(\operatorname{GL}(n,\mathbb R)\) 中，\(X^L_A=AX\)，因此积分曲线满足 \(C'=CX\)、\(C(0)=I\)。矩阵级数
	\[
		e^{tX}=\sum_{k=0}^{\infty}\frac{t^kX^k}{k!}
	\]
	在每个有界时间区间绝对一致收敛，其逐项导数级数也如此，因为可分别用标量指数级数控制矩阵范数。故它可逐项求导，满足所需方程及初值。\(e^{tX}e^{-tX}=I\) 由乘积求导及零点值验证，所以该曲线始终可逆；唯一性给出 \(\exp_{\operatorname{GL}(n,\mathbb R)}X=e^X\)。复矩阵按实变量处理时完全相同。对抽象 Lie 群，\(X\) 未必是矩阵，定义仍是前面的积分曲线，不能预先写一个不存在的矩阵幂级数。
\end{example}

在加法群 \(\mathbb R^n\) 中，指数就是 \(X\mapsto X\)；在正实数乘法群中，\(X^L_x=xX\)，指数就是通常的 \(X\mapsto e^X\)。在圆周群中，把 \(\mathrm{T}_1S^1\) 写成 \(i\mathbb R\)，便有 \(\exp_{S^1}(it)=e^{it}\)。这三个例子分别显示指数可能是恒等映射、一般微分同胚，也可能因周期而失去单射性。

\subsection{指数映射的微分和局部逆}
\label{g1:subsec:1602:02}

\begin{proposition}[指数坐标]\label{g1:prop:exponential-local-diffeomorphism}
	设 \(G\) 为有限维实 Lie 群。用向量空间的平移识别 \(\mathrm{T}_0\mathfrak g=\mathfrak g\)，则
	\[
		(\mathrm{d}\exp_G)_0=\operatorname{id}_{\mathfrak g}.
	\]
	因而零点有一个开邻域 \(V\subset\mathfrak g\)，单位元有一个开邻域 \(U\subset G\)，使 \(\exp_G:V\to U\) 为微分同胚。其局部逆记为 \(\log:U\to V\)，称由此得到的坐标为指数坐标。
\end{proposition}
\begin{proof}
	对 \(Y\in\mathfrak g\)，曲线 \(t\mapsto tY\) 在零点速度为 \(Y\)，而
	\[
		(\mathrm{d}\exp_G)_0Y
		=\left.\frac{\mathrm{d}}{\mathrm{d}t}\right|_{t=0}\exp_G(tY)
		=c_Y'(0)=Y.
	\]
	于是微分可逆。\cref{g1:thm:inverse-function}适用于这两个同维无边界流形间的光滑映射，给出所需邻域及光滑局部逆。必要时可继续缩小 \(V\)，使其关于原点对称且为星形。
\end{proof}

结论只保证单位元附近的可逆性。圆周的指数已有 \(2\pi\) 周期；一般连通群的指数甚至未必满射。不过，单位元的小邻域已足以生成整个单位连通分支。

\begin{proposition}[由指数生成单位分支]\label{g1:prop:exponentials-generate-identity-component}
	设 \(G\) 为有限维实 Lie 群，\(G^0\) 为包含单位元的连通分支。则 \(G^0\) 是开正规子群，并且由任意足够小的对称指数邻域 \(U=\exp_G(V)\) 生成。特别地，连通 Lie 群的每个元素都是有限个指数的乘积。
\end{proposition}
\begin{proof}
	流形局部道路连通，故道路分支均开，且连通分支就是道路分支。若 \(g,h\) 都能用路径从 \(e\) 到达，路径的逐点乘积和逐点逆给出到 \(gh\)、\(g^{-1}\) 的路径，所以 \(G^0\) 为子群。共轭把单位元保持不变，又把路径变成路径，故 \(G^0\) 正规且开。

	取包含于 \(G^0\) 的对称指数邻域 \(U\)。令 \(K=\bigcup_{m\ge1}U^m\)。由于 \(e\in U=U^{-1}\)，这确是由 \(U\) 生成的子群。它开，因为每个有限乘积集是开集的平移之并；它在 \(G^0\) 中的每个陪集也开，因此补集开。于是 \(K\) 在连通空间 \(G^0\) 中既开又闭且非空，只能等于 \(G^0\)。
\end{proof}

为了计算局部乘法的高阶项，还需要知道非零点处的指数微分。这里应先把输出切向量移回单位元，才能与输入 \(Y\in\mathfrak g\) 比较。定义一个 \(\mathfrak g\) 值一形式
\[
	\theta_g=(\mathrm{d}L_{g^{-1}})_g:\mathrm{T}_gG\longrightarrow\mathfrak g.
\]
它满足 \(\theta(X^L)=X\)。在矩阵群中，\(\theta_g(v)=g^{-1}v\)；一般 Lie 群没有矩阵乘法，也仍由左平移微分定义这个对象。

\begin{proposition}[指数微分的积分公式]\label{g1:prop:lie-exponential-differential}
	设 \(G\) 为有限维实 Lie 群，\(\mathfrak g=\mathrm{T}_eG\)，\(X,Y\in\mathfrak g\)。令 \(\theta_g=(\mathrm{d}L_{g^{-1}})_g:\mathrm{T}_gG\to\mathfrak g\)，并记 \(\operatorname{ad}_X(Y)=[X,Y]\)，则
	\begin{equation}\label{g1:eq:lie-dexp-integral}
		\theta_{\exp_GX}\bigl((\mathrm{d}\exp_G)_XY\bigr)
		=J_XY,\qquad
		J_X=\int_0^1 e^{-s\operatorname{ad}_X}\,\mathrm{d}s.
	\end{equation}
	右侧的指数是有限维向量空间 \(\mathfrak g\) 上线性算子的指数，因而有收敛展开
	\begin{equation}\label{g1:eq:lie-dexp-series}
		J_X=\sum_{k=0}^{\infty}\frac{(-1)^k}{(k+1)!}\operatorname{ad}_X^k
		=I-\tfrac12\operatorname{ad}_X+\tfrac16\operatorname{ad}_X^2+O(\|X\|^3).
	\end{equation}
\end{proposition}
\begin{proof}
	先计算 \(\theta\) 的外微分。将\eqref{g1:eq:exterior-one-form-evaluation}逐分量应用于 \(\mathfrak g\) 值形式，对左不变场有
	\[
		\mathrm{d}\theta(X^L,Y^L)=-\theta([X^L,Y^L])=-[X,Y].
	\]
	两边都是对切向输入逐点双线性的二形式，而左不变场的点值张成每个切空间，故一般输入 \(V,W\) 满足 \(\mathrm{d}\theta(V,W)=-[\theta(V),\theta(W)]\)。

	取二参数映射 \(a(s,t)=\exp_G(s(X+tY))\)，并记
	\[
		v(s,t)=\theta(\partial_sa),\qquad u(s,t)=\theta(\partial_ta).
	\]
	对 \(a^*\theta=v\,\mathrm{d}s+u\,\mathrm{d}t\) 作外微分。外微分与拉回交换，见\cref{g1:prop:exterior-pullback}，所以刚才的恒等式给出
	\[
		\partial_su-\partial_tv=-[v,u].
	\]
	固定 \(t\) 时，\(s\mapsto a(s,t)\) 是初速度 \(X+tY\) 的单参数子群，因此 \(v=X+tY\)。在 \(t=0\) 处得到有限维线性方程
	\[
		\partial_su(s,0)=Y-\operatorname{ad}_Xu(s,0),\qquad u(0,0)=0.
	\]
	其解为 \(u(s,0)=\int_0^s e^{-(s-r)\operatorname{ad}_X}Y\,\mathrm{d}r\)：直接求导检查方程与初值，再用唯一性即可确认。令 \(s=1\)，左端是所需的左平移后微分，换元得到\eqref{g1:eq:lie-dexp-integral}。算子指数级数在 \(s\in[0,1]\) 上绝对一致收敛，可以逐项积分；双线性括号在有限维空间中有界，故 \(\|\operatorname{ad}_X\|\le C\|X\|\)，也就得到余项估计。
\end{proof}

公式中分母 \((k+1)!\) 使级数对每个 \(X\) 收敛，但 \(J_X\) 未必处处可逆。指数的局部可逆性在零点附近来自 \(J_0=I\)；离开这个邻域以后，不能仅凭级数存在就保证没有奇点。

\subsection{共轭与指数的相容}
\label{g1:subsec:1602:03}

\begin{proposition}[指数的自然性]\label{g1:prop:exponential-naturality}
	设 \(F:G\to H\) 为有限维实 Lie 群的光滑同态，\(\varphi=(\mathrm{d}F)_{e_G}\)。则对每个 \(X\in\mathfrak g\)，有
	\begin{equation}\label{g1:eq:exponential-homomorphism}
		F(\exp_GX)=\exp_H(\varphi X).
	\end{equation}
	特别地，若 \(g\in G\)，则
	\begin{equation}\label{g1:eq:exponential-conjugation}
		g\exp_G(X)g^{-1}=\exp_G(\operatorname{Ad}_gX),\qquad
		\operatorname{Ad}_{\exp_GX}=e^{\operatorname{ad}_X}.
	\end{equation}
\end{proposition}
\begin{proof}
	复合曲线 \(t\mapsto F(\exp_G(tX))\) 是 \(H\) 中的单参数子群，零点速度为 \(\varphi X\)。单参数子群的唯一性因此使它等于 \(t\mapsto\exp_H(t\varphi X)\)，令 \(t=1\) 得第一式。取同态 \(F=C_g\)、\(C_g(h)=ghg^{-1}\)，其单位元微分是 \(\operatorname{Ad}_g\)，便得共轭公式。

	再取 \(F=\operatorname{Ad}:G\to\operatorname{GL}(\mathfrak g)\)。\cref{g1:prop:adjoint-representation}及\eqref{g1:eq:ad-bracket}说明它是光滑同态，单位元微分为 \(X\mapsto\operatorname{ad}_X\)。刚证明的自然性给出 \(\operatorname{Ad}_{\exp_GX}=\exp_{\operatorname{GL}(\mathfrak g)}(\operatorname{ad}_X)\)，而\cref{g1:exm:matrix-exponential-series}把右端识别为算子指数。
\end{proof}

两个公式分别联系群上的共轭、单位切空间上的线性作用及 Lie 代数的括号。在矩阵群中，它们成为 \(g e^Xg^{-1}=e^{gXg^{-1}}\) 和 \(e^XYe^{-X}=e^{\operatorname{ad}_X}Y\)。后一式的前几项是 \(Y+[X,Y]+\tfrac12[X,[X,Y]]+\cdots\)，因此括号描述共轭造成的最初变化。

\begin{corollary}[交换方向的指数]\label{g1:cor:commuting-exponentials}
	设 \(G\) 为有限维实 Lie 群，\(X,Y\in\mathfrak g\)，且 \([X,Y]=0\)。则对任意实数 \(s,t\)，\(\exp_G(sX)\) 与 \(\exp_G(tY)\) 交换，并且 \(\exp_G(X+Y)=\exp_G(X)\exp_G(Y)\)。
\end{corollary}
\begin{proof}
	由于 \(\operatorname{ad}_XY=0\)，有 \(e^{s\operatorname{ad}_X}Y=Y\)。将其代入共轭公式，得到两条单参数子群的所有元素交换。令 \(c(t)=\exp_G(tX)\exp_G(tY)\)。因这些元素交换，\(c(s+t)=c(s)c(t)\)，故 \(c\) 是单参数子群。群乘法在 \((e,e)\) 处的微分为 \((X,Y)\mapsto X+Y\)：分别固定一个因子时都是恒等映射，微分又线性。因而 \(c'(0)=X+Y\)，唯一性使 \(c(t)=\exp_G(t(X+Y))\)，令 \(t=1\) 即得结论。
\end{proof}

\begin{corollary}[连通源上的唯一性]\label{g1:cor:lie-homomorphism-uniqueness}
	设 \(G,H\) 为有限维实 Lie 群，且 \(G\) 连通。若两个光滑同态 \(F_1,F_2:G\to H\) 在单位元处具有相同微分，则 \(F_1=F_2\)。
\end{corollary}
\begin{proof}
	自然性说明两同态在全部指数上相同。由\cref{g1:prop:exponentials-generate-identity-component}，\(G\) 的每个元素都是有限个指数的乘积；同态保持乘积，所以在每个群元素上都相同。
\end{proof}

这个结论提供唯一性，没有提供存在性。给定一个 Lie 代数同态能否成为群同态的微分，还可能受到环路的限制；\secref{g1:sec:1604}将具体构造延拓并解释单连通性的作用。若源群不连通，单位元微分甚至不记录其他连通分支上的离散选择，例如离散群的单位切空间总是零。

\subsection{局部乘法与 BCH 的三阶展开}
\label{g1:subsec:1602:04}

缩小指数坐标域，使 \(X,Y\) 充分小时 \(\exp_GX\exp_GY\) 仍在 \(\log\) 的定义域，记
\[
	\mu(X,Y)=\log(\exp_GX\exp_GY).
\]
这个局部运算把原来的群乘法写回向量空间。若 \(X,Y\) 交换，上一推论说明它就是 \(X+Y\)；不交换时，括号给出修正。下面证明\term[BCHgongshi]{Baker--Campbell--Hausdorff 公式}[Baker--Campbell--Hausdorff formula]的三阶形式，简称 BCH 公式。\footnote{Baker，Henry Frederick Baker（1866--1956），英国数学家，研究几何与分析；Campbell，John Edward Campbell（1862--1924），生于爱尔兰、在英国任教的数学家，研究连续变换群与微分几何。Hausdorff 的人物介绍见\chapref{g1:ch:04}。}

\begin{proposition}[BCH 的三阶展开]\label{g1:prop:bch-finite-order}
	设 \(G\) 为有限维实 Lie 群，在 \(\mathfrak g\) 上任选一个范数。对充分小的 \(X,Y\)，有
	\begin{equation}\label{g1:eq:bch-third-order}
		\begin{aligned}
			\log(\exp_GX\exp_GY)
			={}&X+Y+\tfrac12[X,Y]\\
			&+\tfrac1{12}[X,[X,Y]]+\tfrac1{12}[Y,[Y,X]]
			+O((\|X\|+\|Y\|)^4).
		\end{aligned}
	\end{equation}
	余项是在某个共同邻域内、相对于所选范数的统一估计。本命题只陈述有限阶 Taylor 展开。
\end{proposition}
\begin{proof}
	我们把第二个因子逐渐乘进去，比较这条运动的左平移后速度。对 \(0\le t\le1\)，令
	\[
		Z(t)=\log(\exp_GX\exp_G(tY)).
	\]
	只要 \(X,Y\) 在足够小的共同邻域中，整条路径都在 \(\log\) 的定义域。其左平移后的速度是 \(Y\)：乘上固定的左因子 \(\exp_GX\) 不改变 \(\theta\) 对速度的值，而 \(\exp_G(tY)\) 的该值本来就是 \(Y\)。由\eqref{g1:eq:lie-dexp-integral}，得到
	\[
		J_{Z(t)}Z'(t)=Y,\qquad Z(0)=X.
	\]
	因为 \(J_0=I\)，小邻域内 \(J_Z\) 可逆。将\eqref{g1:eq:lie-dexp-series}乘以待定的二次多项式，逐次比较可得
	\[
		J_Z^{-1}=I+\tfrac12\operatorname{ad}_Z+\tfrac1{12}\operatorname{ad}_Z^2+O(\|Z\|^3).
	\]
	确切地说，右侧多项式与 \(J_Z\) 的乘积是 \(I+O(\|Z\|^3)\)；再乘以在小邻域上一致有界的 \(J_Z^{-1}\)，便证明所列余项。因此
	\begin{equation}\label{g1:eq:bch-evolution}
		Z'=Y+\tfrac12[Z,Y]+\tfrac1{12}[Z,[Z,Y]]+O(\delta^4),
		\qquad \delta=\|X\|+\|Y\|.
	\end{equation}
	这里 \(Z(t)=O(\delta)\) 对 \(0\le t\le1\) 一致成立：映射 \((t,X,Y)\mapsto Z(t)\) 光滑，在 \((X,Y)=(0,0)\) 为零，沿连接零点的线段积分其一阶微分，并利用紧时间区间上的共同界即可。故逆算子余项作用于 \(Y\) 后确为 \(O(\delta^4)\)。

	先从上式积分，得到 \(Z=X+tY+O(\delta^2)\)。再代入一次，利用 \([Y,Y]=0\)，得到
	\[
		Z=X+tY+\tfrac t2[X,Y]+O(\delta^3).
	\]
	将这个精度代入一次括号项，将 \(X+tY+O(\delta^2)\) 代入二重括号项，便有
	\begin{align*}
		Z'
		&=Y+\tfrac12[X,Y]+\tfrac t4[[X,Y],Y]\\
		&\quad+\tfrac1{12}[X+tY,[X,Y]]+O(\delta^4)\\
		&=Y+\tfrac12[X,Y]+\tfrac1{12}[X,[X,Y]]
		-\tfrac t6[Y,[X,Y]]+O(\delta^4).
	\end{align*}
	所有余项在同一紧时间区间上一致，积分保持其阶数。从 \(t=0\) 积分到 \(1\)，并使用 \(-[Y,[X,Y]]=[Y,[Y,X]]\)，即得\eqref{g1:eq:bch-third-order}。
\end{proof}

三阶公式的推导不要求把 \(G\) 嵌入矩阵群，也没有用一个未证明收敛的无穷括号级数替代群乘法。对后面闭子群定理的局部论证，二阶及其余项控制已经够用；保留三阶则能看清两个方向在高一阶怎样再次相互影响。

\begin{corollary}[乘积极限]\label{g1:cor:lie-product-limit}
	设 \(G\) 为有限维实 Lie 群，\(X,Y\in\mathfrak g\)、\(t\in\mathbb R\)。则
	\[
		\left(\exp_G\left(\frac{tX}{n}\right)
		\exp_G\left(\frac{tY}{n}\right)\right)^n
		\longrightarrow\exp_G(t(X+Y))\qquad(n\to\infty).
	\]
\end{corollary}
\begin{proof}
	当 \(n\) 充分大时，小乘积在对数邻域中。令其对数为 \(Z_n\)，BCH 的二阶精度给出
	\[
		Z_n=\frac t n(X+Y)+O(n^{-2}),\qquad nZ_n\longrightarrow t(X+Y).
	\]
	每个小乘积恰等于 \(\exp_GZ_n\)，而单参数子群的乘法规律给出 \((\exp_GZ_n)^n=\exp_G(nZ_n)\)。最后使用指数的连续性即可。这里取幂的是同一个群元素，无须交换两个不同方向的指数。
\end{proof}

\begin{example}\label{g1:exm:heisenberg-bch}
	三阶上三角幺幂矩阵是对角条目均为 \(1\) 的上三角矩阵；三个上三角条目给出 \(\mathbb R^3\) 坐标，乘法和求逆为多项式，因此它们构成 Lie 群。在其 Lie 代数中取 \(X=aE_{12}+bE_{23}+cE_{13}\)，其中 \(E_{ij}\) 表示只有第 \((i,j)\) 项为 \(1\) 的矩阵。因 \(X^2=abE_{13}\)、\(X^3=0\)，有
	\[
		\exp X=I+aE_{12}+bE_{23}+\left(c+\tfrac12ab\right)E_{13}.
	\]
	对 \(Y=a'E_{12}+b'E_{23}+c'E_{13}\)，有 \([X,Y]=(ab'-a'b)E_{13}\)，且所有二重括号都为零。直接将上面的两个指数矩阵相乘，再用同一公式读回 \((a+a',b+b',c+c'+\tfrac12(ab'-a'b))\)，得到精确等式
	\[
		\log(\exp X\exp Y)=X+Y+\tfrac12[X,Y].
	\]
	精确性来自本例的实际矩阵计算；一般 Lie 群不能只因有限阶展开里暂时没有写出高阶项，就把余项删去。
\end{example}

% !TeX root = ../../Book1.tex
% 纲要标题：### 16.3 Lie 子代数的积分与闭子群定理
% 纲要位置：几何学计划纲要.md，第 751 行。
% 正文状态：初稿；尚未正式审核。

% 纲要内容开始
% 
% * Lie 子代数的积分分布
% * 浸入子群的构造
% * 闭子群定理
% * 闭性与嵌入性的关系
% 
% 纲要内容结束

\section{Lie 子代数的积分与闭子群定理}
\label{g1:sec:1603}

单位元处的一组允许方向，怎样长成一个子群？线性子空间还不够，因为沿两种方向交替运动会产生括号方向；要求这些新方向仍被允许，就得到 Lie 子代数。Frobenius 定理将把它积分成经过单位元的最大叶，但这片叶首先只有自己的流形拓扑，未必采用原群的子空间拓扑。闭子群定理将进一步说明，若一个代数子群在原群中闭，单凭这个拓扑条件就足以恢复嵌入子流形结构。

\ProseParagraphBreak

以下 Lie 群均为有限维实、Hausdorff、第二可数且无边界。固定 Lie 群 \(G\)，单位元为 \(e\)，Lie 代数为 \(\mathfrak g=\mathrm{T}_eG\)。我们使用\cref{g1:thm:frobenius,g1:thm:maximal-connected-leaf}中已经构造的最大积分叶，并使用上一节的指数映射及局部逆。本节只讨论子群自身的流形结构；陪集空间的流形结构将在下一章构造。

\subsection{Lie 子代数的积分分布}
\label{g1:subsec:1603:01}

回顾前面的定义：Lie 子代数是满足 \([\mathfrak h,\mathfrak h]\subseteq\mathfrak h\) 的线性子空间 \(\mathfrak h\subseteq\mathfrak g\)。这个条件可以在整个群上展开：把 \(\mathfrak h\) 从单位元左移到每个点，得到
\begin{equation}\label{g1:eq:subalgebra-left-distribution}
	D_g=(\mathrm{d}L_g)_e\mathfrak h,\qquad g\in G.
\end{equation}
这里 \(L_g(a)=ga\)，所以每个 \(D_g\) 都有相同维数。我们将证明括号封闭性正好保证这个分布可积。

\begin{proposition}[左不变积分分布]\label{g1:prop:lie-subalgebra-left-distribution}
	设 \(G\) 为有限维实 Lie 群，\(\mathfrak g=\mathrm{T}_eG\)，\(\mathfrak h\subseteq\mathfrak g\) 为 \(r\) 维 Lie 子代数。令 \(D_g=(\mathrm{d}L_g)_e\mathfrak h\)，即\eqref{g1:eq:subalgebra-left-distribution}中的分布。则 \(D\) 是秩 \(r\) 的光滑对合分布，所有左平移都保持此分布。对 \(X\in\mathfrak h\)，还有
	\[
		\operatorname{Ad}_{\exp(tX)}\mathfrak h=\mathfrak h\qquad(t\in\mathbb R).
	\]
\end{proposition}
\begin{proof}
	选 \(\mathfrak h\) 的基 \(X_1,\ldots,X_r\)，将其延成左不变向量场 \(X_1^L,\ldots,X_r^L\)。这些场光滑，且在每点线性无关并张成 \(D_g\)，故给出光滑分布。Lie 代数括号的定义及 \(\mathfrak h\) 的封闭性说明，每个 \([X_i^L,X_j^L]\) 都是这些基场的常系数线性组合。对一般局部截面 \(U=\sum_i f_iX_i^L\)、\(V=\sum_j h_jX_j^L\)，有
	\[
		[U,V]=\sum_{i,j}\left(f_ih_j[X_i^L,X_j^L]
		+f_iX_i^L(h_j)X_j^L-h_jX_j^L(f_i)X_i^L\right),
	\]
	仍取值于 \(D\)，所以分布对合。又因为 \(L_aL_g=L_{ag}\)，有 \(\mathrm{d}L_a(D_g)=D_{ag}\)，证明左平移不变性。

	最后固定 \(X,Y\in\mathfrak h\)。由\cref{g1:prop:exponential-naturality}，伴随作用与算子指数满足
	\[
		\operatorname{Ad}_{\exp(tX)}Y
		=e^{t\operatorname{ad}_X}Y
		=\sum_{j=0}^{\infty}\frac{t^j}{j!}\operatorname{ad}_X^jY.
	\]
	括号封闭性说明 \(\operatorname{ad}_X=[X,\cdot]\) 保持 \(\mathfrak h\)，所以级数的每个部分和均在 \(\mathfrak h\) 中。有限维线性子空间闭，级数极限也在其中。因此 \(\operatorname{Ad}_{\exp(tX)}\mathfrak h\subseteq\mathfrak h\)；把 \(t\) 换成 \(-t\)，再用两伴随算子互逆，得到反向包含。
\end{proof}

Frobenius 定理现在给出叶状坐标盒，在其中 \(D\) 由前 \(r\) 个坐标方向张成。后面的横向坐标在每张叶片上恒定。我们会反复使用这个局部事实来判断一个环境光滑映射何时也对叶自身的拓扑光滑。

\begin{lemma}[取值于叶的光滑映射]\label{g1:lem:leaf-valued-map}
	设 \(M,N\) 为有限维、Hausdorff、第二可数的无边界光滑流形，\(D\) 为 \(M\) 上的常秩光滑对合分布。若 \(N\) 连通，\(F:N\to M\) 光滑，且
	\[
		\mathrm{d}F_q(\mathrm{T}_qN)\subseteq D_{F(q)}
	\]
	对每个 \(q\in N\) 成立，则 \(F(N)\) 包含在一片最大叶中。将值域限制为该叶后，\(F\) 对叶的光滑结构仍然光滑。
\end{lemma}
\begin{proof}
	固定 \(q\in N\)，取连通的小坐标邻域 \(V\ni q\)，使 \(F(V)\) 落在一个叶状坐标盒中，盒内坐标记为 \((u,v)\)，其中 \(v\) 是横向坐标。由于 \(\mathrm{d}v\) 在 \(D\) 上为零，链式法则给出 \(\mathrm{d}(v\circ F)=0\)。在连通坐标小球内沿线段积分导数，便知 \(v\circ F\) 恒为一个常值。因此 \(F(V)\) 位于一张叶片中。

	在 \(N\) 中，把像位于同一最大叶的点归成一类。上面的局部结论说明这些类都是开集，它们两两不交并覆盖 \(N\)，故连通性使只有一个非空类。这证明全像位于一片叶。每个 \(F(V)\) 又位于该叶的一张开叶片中，叶片坐标下的 \(F\) 正是光滑函数 \(u\circ F\)，所以到叶的映射局部光滑，进而整体光滑。这里没有要求 \(F\) 是浸入，也没有把叶拓扑替换成环境子空间拓扑。
\end{proof}

\subsection{浸入子群的构造}
\label{g1:subsec:1603:02}

先说清要构造哪一种子群。一个子群 \(H\subseteq G\) 可以带有自己的 Lie 群结构；只要求包含映射光滑且微分单射时，得到的是浸入子群。如果还要求这份拓扑等于 \(G\) 在 \(H\) 上诱导的子空间拓扑，才得到嵌入子群。两种情形中的群乘法都必须对所用的流形结构光滑，不能从集合上的子群性质直接略过这一步。

\begin{definition}[浸入 Lie 子群]\label{g1:def:immersed-lie-subgroup}
	设 \(G\) 为有限维实 Lie 群。子群 \(H\subseteq G\) 若带有有限维、Hausdorff、第二可数的无边界 Lie 群结构，且包含映射 \(i:H\hookrightarrow G\) 为光滑浸入，则称为\term[jinruliziqun]{浸入 Lie 子群}[immersed Lie subgroup]。若包含映射还是光滑嵌入，则称为嵌入 Lie 子群。通过 \(\mathrm{d}i_e\) 将 \(\mathrm{T}_eH\) 识别为 \(\mathfrak g\) 的子空间。
\end{definition}

\begin{theorem}[Lie 子代数的积分]\label{g1:thm:lie-subalgebra-integration}
	设 \(G\) 为有限维实、Hausdorff、第二可数的无边界 Lie 群，\(\mathfrak h\subseteq\mathfrak g\) 为 Lie 子代数。存在唯一的连通浸入 Lie 子群 \(H\subseteq G\)，满足 \(\operatorname{Lie}(H)=\mathfrak h\)。其底集合及流形结构就是左不变分布 \(D_g=(\mathrm{d}L_g)_e\mathfrak h\) 经过 \(e\) 的最大叶，并且
	\begin{equation}\label{g1:eq:integral-subgroup-generated-exponentials}
		H=\langle\exp_G(\mathfrak h)\rangle.
	\end{equation}
	这里右侧表示由全部 \(\exp_G X\)、\(X\in\mathfrak h\) 生成的代数子群。
\end{theorem}
\begin{proof}
	由\cref{g1:prop:lie-subalgebra-left-distribution,g1:thm:maximal-connected-leaf}，经过 \(e\) 的最大叶 \(H\) 是连通 \(r\) 维流形，包含映射为浸入，且 \(\mathrm{T}_gH\) 在 \(\mathrm{T}_gG\) 中的像是 \(D_g\)。先提醒叶拓扑所需的两个可数与分离条件。它使每张局部叶片为开坐标片，包含映射连续；环境中的两个不交开邻域限制到叶上仍然开，所以 \(H\) Hausdorff。第二可数性则来自叶片构造：在 \(G\) 的可数个叶状坐标盒中，一张叶片与任意另一盒的交只有可数个连通分支，每个分支只落在那一盒的一张叶片中。沿有限相交链逐步加入叶片，每一步至多可数，故经过 \(e\) 的叶只需可数张叶片覆盖。取各片可数坐标基之并，得到叶的可数基。这正是最大叶定理中用于叶拓扑的论证，并非从环境子空间拓扑直接继承。

	左平移 \(L_a\) 保持 \(D\)，所以它把经过 \(e\) 的最大叶送到经过 \(a\) 的最大叶：若后者还有更大的连通积分流形，左移回来就会违背前者的最大性。若 \(a\in H\)，经过 \(a\) 的叶仍是 \(H\)，故 \(aH=H\)。于是 \(a,b\in H\) 蕴含 \(ab\in H\)，而 \(e\in aH\) 蕴含 \(a^{-1}\in H\)。这先证明 \(H\) 在集合上是子群。各 \(L_a|_H\)、\(a\in H\)，由\cref{g1:lem:leaf-valued-map}光滑，逆 \(L_{a^{-1}}|_H\) 同样光滑，所以已经是叶流形上的微分同胚。

	若 \(X\in\mathfrak h\)，曲线 \(t\mapsto\exp(tX)\) 的速度为 \(X^L\)，始终切于 \(D\)。它经过 \(e\)，故由取值于叶的映射引理，整条曲线在 \(H\) 中且对叶拓扑光滑。令 \(H'\) 是由这些曲线的值生成的子群，则 \(H'\subseteq H\)。为证明相等，选 \(\mathfrak h\) 的基 \(X_1,\ldots,X_r\)，考虑
	\[
		f(t^1,\ldots,t^r)=\exp(t^1X_1)\cdots\exp(t^rX_r).
	\]
	其值在 \(H'\) 中。记第 \(i\) 个因子之后的乘积为 \(b_i\)。对 \(t^i\) 求导，再用左平移送回单位元，所得方向为 \(\operatorname{Ad}_{b_i^{-1}}X_i\)。每个 \(b_i\) 都是 \(\exp(\mathfrak h)\) 中元素的乘积，上一命题说明它的伴随作用保持 \(\mathfrak h\)，故 \(\mathrm{d}f\) 的像包含于 \(D\)。取值于叶的映射引理使 \(f:\mathbb R^r\to H\) 光滑。在零点，它的微分把标准基送到 \(X_1,\ldots,X_r\)，所以到 \(\mathrm{T}_eH\) 为同构。逆函数定理说明 \(H'\) 包含 \(H\) 中单位元的一个开邻域。

	由 \(H\) 中的左平移已是微分同胚，\(H'\) 是 \(H\) 中的开子群；它的其他左陪集也开，故补集开。连通性迫使 \(H'=H\)，证明\eqref{g1:eq:integral-subgroup-generated-exponentials}。同时，生成元的伴随作用保持 \(\mathfrak h\)，于是
	\[
		\operatorname{Ad}_a\mathfrak h=\mathfrak h\qquad(a\in H).
	\]

	现在验证群运算在叶拓扑下光滑。对 \(a,b\in H\)、\(X,Y\in\mathfrak h\)，把切向量写成左平移的形式，乘法 \(\mu\) 的微分为
	\begin{equation}\label{g1:eq:subgroup-multiplication-differential}
		\mathrm{d}\mu_{(a,b)}\bigl(\mathrm{d}L_aX,\mathrm{d}L_bY\bigr)
		=\mathrm{d}L_{ab}\bigl(\operatorname{Ad}_{b^{-1}}X+Y\bigr)\in D_{ab}.
	\end{equation}
	这个公式可对曲线
	\[
		a\exp(tX)b\exp(tY)
		=ab\,\exp\bigl(t\operatorname{Ad}_{b^{-1}}X\bigr)\exp(tY)
	\]
	在零点求导得到。环境乘法组成的映射 \(H\times H\to G\) 光滑，微分又取值于 \(D\)，所以取值于叶的映射引理表明它作为 \(H\times H\to H\) 也光滑。对逆映射 \(\iota(a)=a^{-1}\)，同理有
	\[
		\mathrm{d}\iota_a(\mathrm{d}L_aX)
		=\mathrm{d}L_{a^{-1}}(-\operatorname{Ad}_aX)\in D_{a^{-1}},
	\]
	这来自 \((a\exp(tX))^{-1}=a^{-1}\exp(-t\operatorname{Ad}_aX)\)。因此逆映射也对叶结构光滑，\(H\) 确实成为 Lie 群。

	包含映射是光滑群同态，单位元处微分像为 \(\mathfrak h\)，所以其 Lie 代数恰为给定子代数。若 \(K\subseteq G\) 是另一个连通浸入 Lie 子群，且单位元切空间同样为 \(\mathfrak h\)，同态与左平移相容使它在每点的切空间像为 \(D\)。因此 \(K\) 是经过 \(e\) 的连通积分流形。最大叶定理给出包含 \(K\to H\)，并保证它为开嵌入。其像又是 \(H\) 的子群，故是开子群；由 \(H\) 连通，像等于 \(H\)。这个开嵌入遂为微分同胚，证明底集合和 Lie 群结构都唯一。若 \(r=0\)，最大叶只是 \(\{e\}\)，上述结论同样成立。
\end{proof}

浸入子群使用叶拓扑，群运算的光滑性由上面的三个性质保证：微分切于分布、横向坐标恒定，以及叶片本身是叶中的开集。对嵌入子群，则可直接调用\cref{g1:prop:embedded-initial-property}中的限制映射性质。

\subsection{闭子群定理}
\label{g1:subsec:1603:03}

现在反向开始：只给 \(G\) 的一个代数子群 \(K\)，暂时不赋予任何流形结构。若它闭，我们可以用全部留在 \(K\) 内的单参数子群识别候选切空间，再证明附近没有额外的横向部分。闭性将在两次取极限时起决定作用。

\begin{lemma}[小元素的极限方向]\label{g1:lem:closed-subgroup-limit-direction}
	设 \(K\) 是有限维实 Lie 群 \(G\) 的闭子群，给 \(\mathfrak g\) 任选 Euclid 范数。若 \(Z_j\ne0\)、\(Z_j\to0\)、\(\exp Z_j\in K\)，且 \(Z_j/\lVert Z_j\rVert\to Z\)，则
	\[
		\exp(tZ)\in K\qquad\text{对每个 }t\in\mathbb R.
	\]
\end{lemma}
\begin{proof}
	固定 \(t\in\mathbb R\)，令 \(\varepsilon_j=\lVert Z_j\rVert\)，选整数 \(m_j\) 使 \(|m_j\varepsilon_j-t|\le\varepsilon_j\)。因为 \(\varepsilon_j\to0\)，有
	\[
		m_jZ_j=(m_j\varepsilon_j)\frac{Z_j}{\varepsilon_j}\longrightarrow tZ.
	\]
	单参数子群规律给出
	\[
		(\exp Z_j)^{m_j}=\exp(m_jZ_j)\in K.
	\]
	指数映射连续，而 \(K\) 闭，故极限 \(\exp(tZ)\) 也属于 \(K\)。整数 \(m_j\) 可以为负，所以论证同时包含正、负时间；\(t\) 任意，结论成立。
\end{proof}

\begin{theorem}[闭子群定理]\label{g1:thm:closed-subgroup}
	设 \(G\) 为有限维实、Hausdorff、第二可数的无边界 Lie 群，\(K\subseteq G\) 为闭的代数子群。则 \(K\) 具有唯一的嵌入 Lie 子群结构，其 Lie 代数为
	\begin{equation}\label{g1:eq:closed-subgroup-lie-algebra}
		\mathfrak k=\{X\in\mathfrak g:\exp(tX)\in K\text{ 对所有 }t\in\mathbb R\}.
	\end{equation}
	对任意线性补空间 \(\mathfrak g=\mathfrak k\oplus\mathfrak m\)，可以选择零点的邻域 \(V\subseteq\mathfrak k\)、\(W\subseteq\mathfrak m\)，使
	\[
		\Psi:V\times W\longrightarrow U,\qquad \Psi(X,Y)=\exp X\exp Y
	\]
	为到 \(G\) 中单位元开邻域的微分同胚，并且
	\begin{equation}\label{g1:eq:closed-subgroup-local-slice}
		K\cap U=\Psi(V\times\{0\})=\exp V.
	\end{equation}
\end{theorem}
\begin{proof}
	先证明\eqref{g1:eq:closed-subgroup-lie-algebra}确实给出 Lie 子代数。对每个固定 \(t\)，\(X\mapsto\exp(tX)\) 连续，所以 \(\mathfrak k\) 是闭集 \(K\) 的这些原像之交，因而在 \(\mathfrak g\) 中闭。它包含零，并且对任意实数倍封闭，因为 \(t(cX)=(tc)X\)。若 \(X,Y\in\mathfrak k\)，由\cref{g1:cor:lie-product-limit}，对每个实数 \(t\) 有
	\[
		\exp(t(X+Y))
		=\lim_{n\to\infty}\bigl(\exp(tX/n)\exp(tY/n)\bigr)^n.
	\]
	右侧每一项都在子群 \(K\) 中，闭性使极限也在 \(K\) 中。因此 \(X+Y\in\mathfrak k\)，\(\mathfrak k\) 是闭线性子空间。

	对 \(k\in K\)，\cref{g1:prop:exponential-naturality}说明指数与共轭相容，故
	\[
		\exp(t\operatorname{Ad}_kX)=k\exp(tX)k^{-1}\in K
		\qquad(X\in\mathfrak k).
	\]
	于是 \(\operatorname{Ad}_k\mathfrak k=\mathfrak k\)，反向包含使用 \(k^{-1}\)。若 \(X,Y\in\mathfrak k\)，则 \(\exp(sX)\in K\)，所以 \(\operatorname{Ad}_{\exp(sX)}Y\) 始终位于闭线性空间 \(\mathfrak k\)。在零点取差商极限，得到
	\[
		[X,Y]=\left.\frac{\mathrm{d}}{\mathrm{d}s}\right|_{s=0}
		\operatorname{Ad}_{\exp(sX)}Y\in\mathfrak k.
	\]
	这证明其括号封闭性。

	取补空间 \(\mathfrak m\)。我们先证明存在 \(\delta>0\)，使
	\[
		Y\in\mathfrak m,\quad \lVert Y\rVert<\delta,\quad \exp Y\in K
		\quad\Longrightarrow\quad Y=0.
	\]
	若不存在，就可取非零 \(Y_j\in\mathfrak m\)，满足 \(Y_j\to0\)、\(\exp Y_j\in K\)。\(\mathfrak m\) 的单位球面紧，故取子列后 \(Y_j/\lVert Y_j\rVert\to Y\in\mathfrak m\)，且 \(\lVert Y\rVert=1\)。由\cref{g1:lem:closed-subgroup-limit-direction}，\(\exp(tY)\in K\) 对所有 \(t\) 成立，即 \(Y\in\mathfrak k\)。这与 \(\mathfrak k\cap\mathfrak m=\{0\}\) 矛盾。若 \(\mathfrak m=\{0\}\)，这一断言直接成立。

	考虑光滑映射 \(\Psi(X,Y)=\exp X\exp Y\)。\cref{g1:prop:exponential-local-diffeomorphism}给出 \(\mathrm{d}(\exp)_0=\operatorname{id}_{\mathfrak g}\)，群乘法在 \((e,e)\) 的微分是向量相加，故
	\[
		\mathrm{d}\Psi_{(0,0)}(X,Y)=X+Y.
	\]
	直和分解使它可逆。逆函数定理给出零点附近的局部微分同胚；再把定义域缩成乘积 \(V\times W\)，并使 \(W\) 包含于上面的 \(\delta\) 球中。这个乘积的像 \(U\) 仍在 \(G\) 中开。

	若 \(k\in K\cap U\)，唯一写成 \(k=\exp X\exp Y\)，其中 \(X\in V\)、\(Y\in W\)。由于 \(X\in\mathfrak k\)，有 \(\exp X\in K\)，所以
	\[
		\exp Y=(\exp X)^{-1}k\in K.
	\]
	横向排除结论强制 \(Y=0\)。反过来，\(X\in V\) 时 \(\exp X\in K\cap U\)。这就证明了完整的切片等式\eqref{g1:eq:closed-subgroup-local-slice}；它排除了附近所有额外分支，而不仅仅给出一个包含于 \(K\) 的参数片。

	将这张切片用 \(L_k\)、\(k\in K\)，移到每个子群点附近，得到 \(K\) 的嵌入子流形结构。它取环境子空间拓扑，Hausdorff 性和第二可数性从 \(G\) 继承；切片坐标转换由环境坐标转换限制得到，故光滑。群乘法与逆映射在 \(G\) 上光滑，其限制值域在 \(K\) 中，\cref{g1:prop:embedded-initial-property}因而保证它们在这个子流形结构下光滑。所以 \(K\) 是嵌入 Lie 子群。单位元的切片为 \(X\mapsto\exp X\)、\(X\in V\)，其微分为包含 \(\mathfrak k\hookrightarrow\mathfrak g\)，故 Lie 代数恰为 \(\mathfrak k\)。最后，嵌入子流形的光滑结构由包含映射的初始性质唯一决定，证明唯一性。
\end{proof}

当 \(\mathfrak k=\{0\}\) 时，局部切片只是单位元，左移以后每个点都孤立，所以 \(K\) 是离散子群。当 \(\mathfrak k=\mathfrak g\) 时，切片占满单位元的一个开邻域，所以 \(K\) 是开子群。两种极端都包含在同一证明中。

\ProseParagraphBreak

Lie 代数只描述单位元所在的连通部分，不能确定全部离散分支。例如加法群 \(\mathbb R\) 中，\(\{0\}\) 与 \(\mathbb Z\) 都是闭子群，二者 Lie 代数都为零，但它们不是同一个子群。因此前面的积分唯一性明确要求连通。对一般闭子群 \(K\)，单位元连通分支 \(K^0\) 是开连通子群：流形局部路径连通使连通分支开，乘法与逆映射又把含单位元的连通集送入同一分支。它具有与 \(K\) 相同的单位元切空间，所以正是 \(\mathfrak k\) 的唯一连通积分子群。

\subsection{闭性与嵌入性的关系}
\label{g1:subsec:1603:04}

一般嵌入子流形不必闭，例如平面中的开线段。然而子群的平移结构会把单位元附近的局部闭性传播到整个群。另一方面，如果一个已有自身拓扑的浸入 Lie 子群碰巧具有闭像，还须说明原拓扑与闭子群定理给出的拓扑一致，才能说原包含映射是嵌入。

\begin{proposition}[闭性与嵌入性]\label{g1:prop:immersed-subgroup-closed-embedded}
	设 \(H\) 为有限维实 Lie 群 \(G\) 的浸入 Lie 子群，两个 Lie 群都为 Hausdorff、第二可数。则包含映射 \(H\hookrightarrow G\) 是嵌入，当且仅当子集 \(H\) 在 \(G\) 中闭。
\end{proposition}
\begin{proof}
	先设包含映射是嵌入。由\cref{g1:prop:embedding-slice-characterization}，存在 \(G\) 中单位元的开邻域 \(U\)，使 \(H\cap U\) 是一张坐标切片，因而在 \(U\) 中闭。令 \(K=\overline H\)。乘法与逆映射连续，故 \(K\) 仍是子群：逼近 \(a,b\in K\) 的子群点列，其乘积及逆分别逼近 \(ab,a^{-1}\)，这些极限仍在闭包中。因为 \(U\) 开，局部闭性给出 \(K\cap U=H\cap U\)。于是 \(H\) 包含 \(K\) 中单位元的开邻域，故是 \(K\) 的开子群。每个左陪集也开，所以 \(H\) 的补集在 \(K\) 中开，即 \(H\) 又在 \(K\) 中闭。它在 \(K=\overline H\) 中稠密，只能有 \(H=K\)，证明闭性。

	反过来，假设像 \(H\) 在 \(G\) 中闭，令 \(K\) 表示同一个集合上由闭子群定理给出的嵌入 Lie 群结构。原包含映射通过 \(K\) 分解为一个双射光滑群同态 \(j:H\to K\)，光滑性使用嵌入子流形的初始性质。需要证明 \(j^{-1}\) 连续，而不是把这一点当作双射的自动结论。

	给定 \(H\) 中单位元的任意开邻域 \(O\)，可取较小开邻域 \(V\)，使 \(C=\overline V\) 紧且 \(C^{-1}C\subseteq O\)。具体可先由乘法与逆的连续性取 \(W\) 满足 \(W^{-1}W\subseteq O\)，再在坐标图中取闭包包含于 \(W\) 的小球作为 \(V\)。第二可数性使覆盖 \(\{hV:h\in H\}\) 有可数子覆盖，故对某列 \(h_\ell\) 有
	\[
		K=\bigcup_{\ell=1}^{\infty}j(h_\ell)\,j(C).
	\]
	每个 \(j(C)\) 都紧，在 Hausdorff 流形 \(K\) 中闭。它不能没有内点。否则右侧是可数个闭的无内点集；在 \(K\) 的一张坐标图中先取一个具有非空内部的紧闭球，再逐次在前一闭球的内部选一个更小的非空闭球，避开第 \(\ell\) 个闭集。这些嵌套紧集有公共点，而该点不属于任何一个所列闭集，与覆盖矛盾。零维时闭的无内点集只能为空，结论直接成立。

	因此可取 \(K\) 中非空开集 \(W_0\subseteq j(C)\) 及 \(k_0\in W_0\)。则
	\[
		k_0^{-1}W_0\subseteq j(C)^{-1}j(C)
		=j(C^{-1}C)\subseteq j(O).
	\]
	左侧是单位元的开邻域，所以 \(j(O)\) 包含单位元邻域。在 \(O\) 的每个点先平移一个更小的单位元邻域，并应用相同论证，可知 \(j(O)\) 本身开。再用群平移，\(j\) 把任意开集送到开集。于是双射 \(j\) 为同胚，原 \(H\) 的拓扑等于 \(G\) 的子空间拓扑。原包含映射又已知是浸入，故按定义为光滑嵌入。
\end{proof}

这个证明中的第二可数性用于取得可数紧集覆盖；紧性和环境的 Hausdorff 性用于把它们的像看成闭集。它们共同排除了闭像还带有一套更细 Lie 群拓扑的可能。对于一般的单射浸入映射，没有群平移结构时，不能照搬这个结论。

\begin{example}[无理斜率子群]\label{g1:exm:irrational-torus-subgroup}
	在环面 Lie 群 \(T^2=\mathbb R^2/\mathbb Z^2\) 中，固定无理数 \(\alpha\)，定义
	\[
		f_\alpha:\mathbb R\longrightarrow T^2,\qquad
		f_\alpha(t)=[t,\alpha t].
	\]
	它是光滑群同态，微分像为直线 \(\mathfrak h_\alpha=\mathbb R(1,\alpha)\)。若 \(f_\alpha(t)=0\)，则 \(t\in\mathbb Z\) 且 \(\alpha t\in\mathbb Z\)，无理性迫使 \(t=0\)，所以它单射；微分处处非零，故为浸入。将像 \(H_\alpha\) 赋予从 \(\mathbb R\) 传来的 Lie 群结构，便得到 \(\mathfrak h_\alpha\) 的连通积分子群。

	\cref{g1:prop:linear-torus-leaves}已经证明这条轨道在环面中稠密。其关键是整数倍 \(\alpha\) 模 \(1\) 稠密：对任意指定的第一坐标 \(u\)，时刻 \(u+n\) 的第二坐标为 \(\alpha u+n\alpha\)，因此能够进入任意给定小矩形。这个像却不是整个环面，因为它与第一坐标为零的圆周只相交于可数个点 \([0,n\alpha]\)。于是 \(H_\alpha\) 是真稠密子群，既不闭也不嵌入。

	拓扑差异还可直接检测。前述整数逼近给出 \(|n_j|\to\infty\)，但 \(n_j\alpha\bmod1\to0\)，所以 \(f_\alpha(n_j)\to f_\alpha(0)\) 在环面中成立；在 \(H_\alpha\) 自身的参数坐标里，\(n_j\) 却不趋于零。这正是从子空间像返回叶参数的逆映射不连续。子群自身始终是一维 Lie 群，闭包却是二维环面，其 Lie 代数也从一条直线增大为整个 \(\mathbb R^2\)。
\end{example}

因此，Lie 子代数的积分定理提供连通浸入子群，闭子群定理提供采用子空间拓扑的嵌入子群，两者处理的条件不同。后面构造齐性空间时，将用闭性把子群附近的横向坐标组织为陪集空间的光滑坐标。

% 纲要条目排为编号小节标题，并进入目录与书签；说明段仍保留于纲要注释中。

% !TeX root = ../../Book1.tex
% 纲要标题：### 16.4 同态积分及覆盖群
% 纲要位置：几何学计划纲要.md，第 758 行。
% 正文状态：初稿；尚未正式审核。

% 纲要内容开始
% 
% * 局部同态的积分构造
% * 沿路径延拓和同伦无关性
% * 单连通源群上的整体同态
% * 连通 Lie 群的万有覆盖群
% * 离散中心商与同 Lie 代数的不同 Lie 群
% 
% ---
% 
% 纲要内容结束

\section{同态积分及覆盖群}
\label{g1:sec:1604}

Lie 群同态在单位元的微分必为 Lie 代数同态，见\cref{g1:prop:lie-homomorphism-differential}。反过来，一个保持括号的线性映射，能否来自某个群同态？局部问题由微分方程的相容性解决；整体问题还要检查，沿不同路径把这个局部映射继续下去，是否会得到相同结果。圆周上“角度乘以 \(\lambda\)”的例子已经提示：局部允许任意实数 \(\lambda\)，绕完一周后的闭合要求却可能排除其中大部分。

\ProseParagraphBreak

本节始终讨论有限维实 Lie 群，底流形无边界、Hausdorff 且第二可数。记 \(e_G,e_H\) 为单位元，\(\mathfrak g=\mathrm{T}_{e_G}G\)、\(\mathfrak h=\mathrm{T}_{e_H}H\) 为相应 Lie 代数。需要的积分工具是\cref{g1:thm:frobenius,g1:lem:local-integral-plaque}：对合分布局部有叶片，且位于同一叶状坐标盒内的连通积分流形只能处于一张叶片中。整体延拓则使用第五章的连续路径与定端点同伦，拼接 \(\alpha*\beta\) 仍然表示先走 \(\alpha\)，后走 \(\beta\)。

\subsection{局部同态的积分构造}
\label{g1:subsec:1604:01}

\begin{definition}
	\label{g1:def:local-lie-homomorphism}
	设 \(G,H\) 为有限维实 Lie 群，\(U\subset G\) 是单位元的开邻域。光滑映射 \(f:U\to H\) 若满足 \(f(e_G)=e_H\)，并且
	\[
		f(xy)=f(x)f(y)\qquad\text{只要 }x,y,xy\in U,
	\]
	则称为从 \(G\) 到 \(H\) 的\term[jubuliquntongtai]{局部 Lie 群同态}[local Lie group homomorphism]。两个这样的映射局部相同，是指它们在某个共同的单位元邻域上相同。
\end{definition}

要从线性映射 \(\varphi:\mathfrak g\to\mathfrak h\) 找到 \(f\)，可以先问它的图像应有什么切空间。若 \(f\) 是同态，固定 \(g\) 后对 \(f(gx)=f(g)f(x)\) 在 \(x=e_G\) 求导，便得到
\[
	\mathrm{d}f_g\circ\mathrm{d}(L_g)_{e_G}
	=\mathrm{d}(L_{f(g)})_{e_H}\circ\varphi.
\]
因此，图像在 \((g,f(g))\) 处的切向量应当把同一个 \(X\in\mathfrak g\) 送成一对左平移向量。这个要求直接给出待积分的分布。

\begin{theorem}[局部同态积分]
	\label{g1:thm:local-lie-homomorphism}
	设 \(G,H\) 为有限维实 Lie 群，\(\varphi:\mathfrak g\to\mathfrak h\) 为 Lie 代数同态。则存在单位元的连通对称开邻域 \(U\subset G\)，以及局部 Lie 群同态 \(f:U\to H\)，满足 \(\mathrm{d}f_{e_G}=\varphi\)。这样的局部同态在单位元附近唯一。
\end{theorem}
\begin{proof}
	在 \(G\times H\) 上定义秩 \(\dim G\) 的光滑分布
	\begin{equation}
		\label{g1:eq:homomorphism-graph-distribution}
		D_{(g,h)}
		=\left\{\left(\mathrm{d}(L_g)_{e_G}X,\,
		\mathrm{d}(L_h)_{e_H}\varphi X\right):X\in\mathfrak g\right\}.
	\end{equation}
	取 \(\mathfrak g\) 的基 \(X_1,\ldots,X_m\)。对应向量场
	\(Z_i=(X_i^L,(\varphi X_i)^L)\)
	处处线性无关，并张成 \(D\)，因为其第一分量已经线性无关。由\cref{g1:prop:left-invariant-fields,g1:def:lie-algebra-group}以及 \(\varphi\) 保括号，
	\[
		[Z_i,Z_j]
		=\left([X_i,X_j]^L,[\varphi X_i,\varphi X_j]^L\right)
		=\left([X_i,X_j]^L,(\varphi[X_i,X_j])^L\right)
	\]
	仍取值于 \(D\)。一般局部截面是这些 \(Z_i\) 的光滑函数线性组合；括号乘法法则将其括号展开为上述括号与 \(Z_i\) 的函数线性组合，所以 \(D\) 对合。

	由 Frobenius 定理，\((e_G,e_H)\) 附近有叶状坐标盒 \(\Omega\)，并有经过该点的中央叶片 \(P\)。投影 \(\pi_G:G\times H\to G\) 的微分在 \(D\) 上为同构。因此由\cref{g1:thm:inverse-function}，缩小中央叶片后，\(\pi_G|_P\) 是到单位元某开邻域 \(U_0\) 的微分同胚。可以同时缩小叶状坐标盒，使它的整个中央叶片就是这一小片。于是
	\[
		P=\{(g,f(g)):g\in U_0\}
	\]
	定义光滑映射 \(f:U_0\to H\)，且 \(f(e_G)=e_H\)。图像切空间等于 \(D\)，给出
	\begin{equation}
		\label{g1:eq:local-homomorphism-differential}
		\mathrm{d}f_g\circ\mathrm{d}(L_g)_{e_G}
		=\mathrm{d}(L_{f(g)})_{e_H}\circ\varphi.
	\end{equation}
	特别地，在单位元得到 \(\mathrm{d}f_{e_G}=\varphi\)。

	还须证明图像尊重乘法。分布 \(D\) 在 \(G\times H\) 的每个左平移下不变。由乘法及 \(f\) 的连续性，可以选连通对称单位元邻域 \(U\subset U_0\)，使 \(U^2\subset U_0\)，并使
	\[
		(ab,f(a)f(b))\in\Omega\qquad(a,b\in U).
	\]
	这里是对全部 \(a,b\in U\) 的共同缩小：映射 \((a,b)\mapsto(ab,f(a)f(b))\) 在 \((e_G,e_G)\) 取值于开集 \(\Omega\)，故其原像包含一个乘积邻域，再在两个因子的交中取小的连通对称邻域即可。

	固定 \(a\in U\)。映射 \(b\mapsto(ab,f(a)f(b))\)、\(b\in U\)，是中央叶片一部分的左平移，因而是连通积分流形，且全部位于 \(\Omega\) 内。它在 \(b=e_G\) 时经过 \((a,f(a))\in P\)。由\cref{g1:lem:local-integral-plaque}，它只能位于同一张中央叶片 \(P\)。投影到第一分量，便得到
	\[
		f(ab)=f(a)f(b)\qquad(a,b\in U).
	\]
	左边在 \(U_0\) 中有定义；将 \(f\) 限制到 \(U\) 后，尤其得到定义中要求的、当 \(a,b,ab\in U\) 时的局部乘法律。

	最后设 \(\widehat f\) 是另一个微分为 \(\varphi\) 的局部同态。对固定 \(g\) 及充分接近单位元的 \(x\)，对 \(\widehat f(gx)=\widehat f(g)\widehat f(x)\) 求导，得到与\eqref{g1:eq:local-homomorphism-differential}相同的等式。因此 \(\widehat f\) 的图像也是 \(D\) 的积分流形。取足够小的连通定义域，使图像落在 \(\Omega\) 中；它经过 \((e_G,e_H)\)，故也位于中央叶片 \(P\)。中央叶片已经是 \(f\) 的图像，所以两映射在这个共同邻域内相同。
\end{proof}

保持括号的条件在这里承担了明确的相容性责任：它使图分布对合，因而存在所需局部图像。局部唯一性则意味着，后面沿路径延拓时不会在每个小邻域重新引入任意选择。此处得到的仍是局部同态；要把定义域扩展到整个群，必须进一步比较路径。

\subsection{沿路径延拓和同伦无关性}
\label{g1:subsec:1604:02}

设 \(f:U\to H\) 是上面得到的局部同态，\(U=U^{-1}\) 为单位元邻域。从点 \(g\) 走到附近的点 \(g'\)，其左增量是 \(g^{-1}g'\)。只要增量足够小，就可以用 \(f(g^{-1}g')\) 记录目标群中的相应变化。沿整条路径把这些小变化相乘，便是局部同态继续到远处的候选值。

\begin{lemma}[局部同态的路径延拓]
	\label{g1:lem:local-hom-path-extension}
	设 \(G,H\) 为有限维实 Lie 群，\(U=U^{-1}\subset G\) 是单位元开邻域，\(f:U\to H\) 为局部同态。对任意连续路径 \(\gamma:[0,1]\to G\)，可以选有限分割 \(0=t_0<\cdots<t_N=1\)，使
	\[
		\gamma(s)^{-1}\gamma(t)\in U
		\quad\text{对同一小区间内的所有 }s,t\text{ 成立}.
	\]
	定义
	\begin{equation}
		\label{g1:eq:local-hom-path-product}
		J_f(\gamma)
		=f(\gamma(t_0)^{-1}\gamma(t_1))\cdots
		f(\gamma(t_{N-1})^{-1}\gamma(t_N)).
	\end{equation}
	则此值与这样的分割无关，并在连续定端点同伦下不变。对另一连续路径 \(\eta:[0,1]\to G\)，在 \(\gamma(1)=\eta(0)\) 时令 \(\gamma*\eta\) 表示先走 \(\gamma\) 后走 \(\eta\) 的拼接，并令 \(\bar\gamma(t)=\gamma(1-t)\)。它还满足
	\[
		J_f(\gamma*\eta)=J_f(\gamma)J_f(\eta),\qquad
		J_f(\bar\gamma)=J_f(\gamma)^{-1},\qquad
		J_f(a\gamma)=J_f(\gamma),
	\]
	其中 \((a\gamma)(t)=a\,\gamma(t)\) 表示固定群元素 \(a\) 对路径的左平移。常路径的值为 \(e_H\)。若两个局部同态在单位元附近相同，则它们给出相同的路径值。
\end{lemma}
\begin{proof}
	先取对称单位元邻域 \(W\)，使 \(W^{-1}W\subset U\)。各左平移 \(aW\) 覆盖 \(G\)，同一个 \(aW\) 内任意两点的相对增量属于 \(W^{-1}W\)。将这个覆盖沿 \(\gamma\) 拉回紧区间，由\cref{g1:lem:lebesgue-number}选择有限细分，使每个小区间的像都位于某个 \(aW\) 中，便得到所需分割。

	在一个已有小区间中增加分点 \(s\)，原增量分解为
	\[
		\gamma(t_{i-1})^{-1}\gamma(t_i)
		=\bigl(\gamma(t_{i-1})^{-1}\gamma(s)\bigr)
		\bigl(\gamma(s)^{-1}\gamma(t_i)\bigr).
	\]
	两个因子及其乘积都在 \(U\) 中，所以局部乘法律使原来的一个 \(f\) 因子等于细分后的两个因子。反复加入分点不改变总乘积。任意两种分割有共同细分，而且细分仍满足整段增量条件，故定义无关。拼接对应接上两组分割，反向对应逆序并逐项取逆；这里 \(f(x^{-1})=f(x)^{-1}\) 来自 \(xx^{-1}=e_G\)。左平移不改变增量，因为 \((a\gamma(s))^{-1}(a\gamma(t))=\gamma(s)^{-1}\gamma(t)\)。这些事实给出所有运算规律。

	现在取连续定端点同伦 \(K:[0,1]^2\to G\)。仍用左平移 \(aW\) 的原像覆盖紧方形，取足够细的有限矩形网格，使每格的像落在某个 \(aW\) 中。若一格四个顶点的像依次为左下 \(a\)、右下 \(b\)、左上 \(c\)、右上 \(d\)，其边路径的 \(J_f\) 值分别是
	\[
		h=f(a^{-1}b),\quad v'=f(b^{-1}d),\quad
		v=f(a^{-1}c),\quad h'=f(c^{-1}d).
	\]
	因为 \(a^{-1}b,b^{-1}d,a^{-1}d\) 都在 \(U\)，以及另一条边路的三个增量也都在 \(U\)，局部乘法律给出
	\begin{equation}
		\label{g1:eq:local-hom-homotopy-cell}
		hv'=f(a^{-1}d)=vh'.
	\end{equation}
	这个等式允许按原来的顺序把经过一格下边、右边的乘积换成左边、上边的乘积。

	具体地，记网格第 \(j\) 行的水平边值为 \(h_{0j},\ldots,h_{m-1,j}\)，竖边值为 \(v_{0j},\ldots,v_{mj}\)，整行乘积为 \(r_j=h_{0j}\cdots h_{m-1,j}\)。逐格应用\eqref{g1:eq:local-hom-homotopy-cell}得到
	\[
		r_jv_{mj}=v_{0j}r_{j+1}.
	\]
	同伦固定端点，所以左右外边的值 \(v_{0j},v_{mj}\) 都为 \(e_H\)，故 \(r_j=r_{j+1}\)。逐行比较便得到两端路径具有相同的 \(J_f\) 值。整个论证没有交换不相邻因子，也没有要求 \(K\) 光滑。

	若 \(f,\widehat f\) 在更小的单位元邻域 \(U'\) 上相同，用 \(U'\) 再作一次上述细分，每个因子都相同。由已经证明的细分独立性，得到 \(J_f=J_{\widehat f}\)。
\end{proof}

这个证明与\cref{g1:thm:flat-path-independence}使用同一种有限网格方法。此处的 \(J_f\) 是按行程顺序相乘的群元素，故拼接公式也是按这个顺序书写；上一章的 \(P_\gamma\) 是作用于向量的线性映射，其函数复合按另一种书写次序进行。记住所比较的对象，便能核对每个公式中的顺序。

\subsection{单连通源群上的整体同态}
\label{g1:subsec:1604:03}

\begin{theorem}[Lie 群同态的积分]
	\label{g1:thm:lie-homomorphism-integration}
	设 \(G,H\) 为有限维实 Lie 群，\(G\) 连通且单连通。每个 Lie 代数同态 \(\varphi:\mathfrak g\to\mathfrak h\) 都唯一地来自一个光滑群同态 \(F:G\to H\)，即 \(\mathrm{d}F_{e_G}=\varphi\)。唯一性只需要源群连通：两个从同一连通 Lie 群出发的光滑群同态，若在单位元有相同微分，则处处相同。
\end{theorem}
\begin{proof}
	由\cref{g1:thm:local-lie-homomorphism}取得局部同态 \(f:U\to H\)。对 \(g\in G\)，取从 \(e_G\) 到 \(g\) 的连续路径 \(\gamma\)，定义 \(F(g)=J_f(\gamma)\)。单连通性包括道路连通性，所以路径存在；同端点的路径定端点同伦，\cref{g1:lem:local-hom-path-extension}因而保证定义与路径无关。常路径给出 \(F(e_G)=e_H\)。

	先证明 \(F\) 光滑。取足够小的单位元坐标球 \(V\subset U\)，使 \(V^{-1}V\subset U\)，并使坐标把 \(e_G\) 送到原点。对 \(x\in V\)，用坐标中的径向线段给出路径 \(\kappa_x:[0,1]\to V\)，从 \(e_G\) 走到 \(x\)。这一整段满足增量条件，所以 \(J_f(\kappa_x)=f(x)\)。固定 \(g\) 及一条从 \(e_G\) 到 \(g\) 的路径，在其后接上 \(g\kappa_x\)，便有
	\[
		F(gx)=F(g)f(x),\qquad x\in V.
	\]
	右边对 \(x\) 光滑，故 \(F\) 在每个 \(gV\) 上光滑。在 \(g=e_G\) 时得到 \(F|_V=f|_V\)，所以 \(\mathrm{d}F_{e_G}=\varphi\)。

	对 \(g,h\in G\)，取路径 \(\alpha:e_G\to g\)、\(\beta:e_G\to h\)。路径 \(\alpha*(g\beta)\) 从 \(e_G\) 到 \(gh\)，由路径值的拼接与左平移规律，
	\[
		F(gh)=J_f(\alpha*(g\beta))
		=J_f(\alpha)J_f(g\beta)=F(g)F(h).
	\]
	因此 \(F\) 为群同态。

	唯一性由\cref{g1:cor:lie-homomorphism-uniqueness}得到。具体地，单位元微分相同使两同态在所有指数上一致，而连通群的每个元素都是有限个指数的乘积，故两同态处处相同。这部分只需要源群连通，不需要再次使用单连通性。
\end{proof}

\begin{corollary}[环路的积分障碍]
	\label{g1:cor:lie-integration-loop-criterion}
	设 \(G,H\) 为有限维实 Lie 群，\(G\) 连通，\(\varphi:\operatorname{Lie}(G)\to\operatorname{Lie}(H)\) 为 Lie 代数同态。取单位元对称开邻域 \(U\subseteq G\) 及光滑局部同态 \(f:U\to H\)，满足 \((\mathrm{d}f)_{e_G}=\varphi\)。对连续路径 \(\gamma:[0,1]\to G\)，取使同段增量 \(\gamma(s)^{-1}\gamma(t)\) 均在 \(U\) 内的有限分割 \(0=t_0<\cdots<t_N=1\)，按时间顺序定义
	\[
		J_f(\gamma)=\prod_{j=1}^N f\bigl(\gamma(t_{j-1})^{-1}\gamma(t_j)\bigr).
	\]
	约定 \(\pi_1(G,e_G)\) 的乘法先走左边的环路，再走右边的环路。则 \(\varphi\) 能积分成 \(G\to H\) 的整体同态，当且仅当
	\[
		J_f(\gamma)=e_H
		\qquad\text{对每条基于 }e_G\text{ 的连续环路 }\gamma\text{ 成立}.
	\]
	映射 \([\gamma]\mapsto J_f(\gamma)\) 本身是 \(\pi_1(G,e_G)\to H\) 的群同态，并只依赖 \(\varphi\)。
\end{corollary}
\begin{proof}
	同伦不变性保证映射在环路类上良定义，拼接式保证它是群同态，局部积分的唯一性则保证构造不依赖局部积分的选择。若已有整体同态 \(F\)，对任意细分使用局部相同及同态性质，路径乘积逐项消去后等于
	\(F(\gamma(0))^{-1}F(\gamma(1))\)，所以环路值为单位元。

	反之，若全部环路值平凡，而 \(\gamma,\eta\) 都从 \(e_G\) 到 \(g\)，则
	\[
		e_H=J_f(\gamma*\bar\eta)=J_f(\gamma)J_f(\eta)^{-1},
	\]
	故两条路径给出相同值。上一证明中 \(F(g)=J_f(\gamma)\) 的构造及其光滑性、乘法性质便仍然适用。
\end{proof}

\begin{example}
	\label{g1:exm:circle-lie-integration-obstruction}
	将 \(S^1=\{e^{it}:t\in\mathbb R\}\) 的 Lie 代数按单位角速度识别为 \(\mathbb R\)。因为括号为零，每个线性映射 \(X\mapsto\lambda X\)、\(\lambda\in\mathbb R\)，都是 Lie 代数同态。其局部积分是在一个小角度区间内定义的
	\[
		f(e^{it})=e^{i\lambda t}.
	\]
	沿正向一周环路细分，小角度增量之和为 \(2\pi\)，所以路径乘积为 \(e^{2\pi i\lambda}\)。由\cref{g1:thm:circle-fundamental}，任意环路由其整数绕数确定；绕数 \(k\) 的路径值为 \(e^{2\pi i\lambda k}\)。因此整体积分存在当且仅当 \(\lambda\in\mathbb Z\)，此时就是 \(F(z)=z^\lambda\)。

	若把源群改成加法群 \(\mathbb R\)，则 \(t\mapsto e^{i\lambda t}\) 对每个实数 \(\lambda\) 都是整体同态。两种源群有同一个一维 Lie 代数；差别是圆周把相差 \(2\pi\) 的角度识别了，目标值也必须尊重这个识别。
\end{example}

\subsection{连通 Lie 群的万有覆盖群}
\label{g1:subsec:1604:04}

对于不是单连通的 \(G\)，可以先把路径的同伦历史保留下来。第五章的道路类万有覆盖已经这样做了；现在还须证明它能够承载群乘法，并且所得结构确实是通常意义下的 Lie 群。

\begin{theorem}[万有覆盖群]
	\label{g1:thm:universal-cover-lie-group}
	设 \(G\) 为连通有限维实 Lie 群。以 \(e_G\) 为基点的道路类万有覆盖
	\(p:(\widetilde G,\widetilde e)\to(G,e_G)\)
	具有以 \(\widetilde e\) 为单位元的唯一光滑 Lie 群结构，使 \(p\) 为光滑覆盖同态，并在每点为局部微分同胚。其核 \(\ker p\) 为离散中心子群，而且自然同构于 \(\pi_1(G,e_G)\)。特别地，连通 Lie 群的基本群为交换群。这个带群结构的覆盖称为\term[wanyoufugaiqun]{万有覆盖群}[universal covering group]。
\end{theorem}
\begin{proof}
	连通 Lie 群道路连通，而坐标小球可缩，所以第五章\cref{g1:thm:path-class-universal-cover}适用。其点为从 \(e_G\) 出发的连续路径的定端点同伦类，投影取终点，总空间已经知道连通且单连通。对路径 \(\alpha,\beta\) 定义逐点乘积 \((\alpha\cdot\beta)(t)=\alpha(t)\beta(t)\)，并令 \(\alpha^\iota(t)=\alpha(t)^{-1}\)。置
	\begin{equation}
		\label{g1:eq:universal-cover-group-product}
		[\alpha][\beta]=[\alpha\cdot\beta],\qquad
		[\alpha]^{-1}=[\alpha^\iota],\qquad
		\widetilde e=[e_{e_G}].
	\end{equation}
	若改变路径代表，逐点相乘或取逆相应的同伦，就得到新路径之间的定端点同伦，所以运算良定义。结合律、单位律和逆元律由逐点的群公理得到，终点映射 \(p\) 也立即保持乘法。此处 \(\alpha^\iota\) 是逐点取群逆，不是把路径参数反向；它仍从单位元出发。

	为了比较第五章的路径拼接，记 \(a=\alpha(1)\)。有
	\begin{equation}
		\label{g1:eq:pointwise-product-concatenation}
		[\alpha\cdot\beta]=[\alpha*(a\beta)].
	\end{equation}
	事实上，映射 \(Q(s,t)=\alpha(s)\beta(t)\) 将方形对角线送成逐点乘积路径，将先沿底边、再沿右边的折线路径送成 \(\alpha*(a\beta)\)。这两条方形内路径的端点都是 \((0,0),(1,1)\)；按坐标作线性插值给出定端点同伦，再复合 \(Q\) 即得所需同伦。\cref{g1:fig:cover-path-product}显示了这两种行进方式。

	接着核对局部光滑结构。取单位元的一个坐标小球 \(U\)，使其坐标像为以零为中心的凸球。记 \(\kappa_x(t)\) 为坐标径向线段从 \(e_G\) 到 \(x\) 的路径。对 \(q=[\alpha]\)、\(a=p(q)\)，
	\[
		x\longmapsto q_x=[\alpha*(a\kappa_x)],\qquad x\in U,
	\]
	正是 \(aU\) 上经过 \(q\) 的一张覆盖片的局部逆，因为在可缩邻域 \(aU\) 内，追加路径的终点决定其路径类。它与\cref{g1:lem:path-class-basis}的拓扑一致。把 \(G\) 的坐标通过 \(p\) 提升到每张片，就得到相容的局部 Euclid 坐标；交叠处的转换由底空间坐标转换给出。

	乘法在这些坐标中光滑，可以直接核算。设 \(r=[\beta]\)、\(b=p(r)\)，并使用 \(q_x,r_y\) 的上述表示。两条拼接路径逐点相乘后，前半段是 \(\alpha\cdot\beta\)，后半段为
	\[
		a\kappa_x(t)b\kappa_y(t)
		=ab\bigl(b^{-1}\kappa_x(t)b\kappa_y(t)\bigr).
	\]
	固定一个足够小的可缩单位元坐标球 \(W\)。缩小 \(x,y\) 的范围，可以使括号中的整条路径都留在 \(W\) 中：关于 \((x,y,t)\) 的表达连续，在 \(x=y=e_G\) 时对全部 \(t\in[0,1]\) 都为单位元，紧区间给出共同的缩小范围。其终点为 \(w=b^{-1}xby\)。在 \(W\) 中，这条路径与径向路径 \(\kappa_w\) 定端点同伦，故
	\[
		q_xr_y=[(\alpha\cdot\beta)*(ab\kappa_w)],\qquad w=b^{-1}xby.
	\]
	因此乘法的局部坐标表达就是 \(w=b^{-1}xby\)，为光滑映射。类似地，对 \(q_x\) 的代表逐点取逆，其末段为
	\[
		(a\kappa_x(t))^{-1}
		=a^{-1}\bigl(a\kappa_x(t)^{-1}a^{-1}\bigr).
	\]
	缩小后括号内路径留在一个可缩小邻域中，取逆的局部坐标表达为 \(x\mapsto ax^{-1}a^{-1}\)，也光滑。这也先行证明了运算关于覆盖拓扑的连续性。

	还要核对 Hausdorff 性和第二可数性。不同投影点可由底空间的不交开集分离；投影相同的不同覆盖点，则位于某均匀覆盖邻域上方的不交片内，所以 \(\widetilde G\) Hausdorff。第二可数性不能仅由底空间第二可数推出。取一个对称开单位元邻域 \(V\subset\widetilde G\)，包含在一张覆盖坐标片内；它本身第二可数。由于运算已经连续，\(\bigcup_{n\geq1}V^n\) 是开子群，其各陪集也开；\(\widetilde G\) 连通，所以这个开子群等于整个 \(\widetilde G\)。每个乘法映射
	\[
		V^{\times n}\longrightarrow V^n,\qquad
		(v_1,\ldots,v_n)\longmapsto v_1\cdots v_n
	\]
	都是开满射：乘积空间中的开矩形映成若干群开集的乘积，而该乘积是开集，因为可固定前面因子再左平移最后一个开集，随后取并。因此 \(V^{\times n}\) 的可数基的像给出 \(V^n\) 的可数基。各 \(V^n\) 在 \(\widetilde G\) 中开，且可数个这样的集合覆盖 \(\widetilde G\)，所有这些基的并便是总空间的可数基。至此，提升的坐标使 \(\widetilde G\) 成为有限维实 Lie 群，\(p\) 是局部微分同胚。

	其核是覆盖的单位元纤维，因而离散。固定 \(k\in\ker p\)，连续映射 \(q\mapsto qkq^{-1}\) 的像仍在 \(\ker p\) 中。连通空间到离散空间的连续映射为常值，在 \(q=\widetilde e\) 处取值为 \(k\)，故 \(k\) 与每个 \(q\) 交换。因此核位于中心。核的点恰是基于 \(e_G\) 的环路类；对两条环路，\eqref{g1:eq:pointwise-product-concatenation}中的 \(a=e_G\)，所以核乘法正是第五章基本群的拼接乘法。于是 \(\ker p\cong\pi_1(G,e_G)\)，而中心的子群交换，得到最后结论。

	最后说明唯一性。使 \(p\) 成为局部微分同胚的光滑结构由各覆盖片唯一确定。若还有另一套使 \(p\) 为同态的连续群乘法，那么两种乘法都是映射 \((q,r)\mapsto p(q)p(r)\) 到覆盖空间的提升，并在 \((\widetilde e,\widetilde e)\) 取相同值。空间 \(\widetilde G\times\widetilde G\) 连通，故\cref{g1:lem:lift-uniqueness}使两种乘法相同；单位元和逆元也随之确定。
\end{proof}

% !TeX root = ../Book1.tex
% 路径逐点相乘与先拼接再左平移的定端点同伦。
\begin{figure}[htbp]
	\centering
	\begin{tikzpicture}[font=\small,scale=1.05]
		\fill[GeometryBlueLight] (0,0) rectangle (3.8,3.2);
		\draw[GeometryGray] (0,0) rectangle (3.8,3.2);
		\draw[GeometryBlue,very thick,->] (0,0)--(3.8,3.2) node[midway,sloped,above] {\(\alpha\cdot\beta\)};
		\draw[orange!75!black,very thick,->] (0,0)--(3.8,0) node[midway,below] {\(\alpha\)};
		\draw[orange!75!black,very thick,->] (3.8,0)--(3.8,3.2) node[midway,right] {\(a\beta\)};
		\fill[GeometryBlue] (0,0) circle[radius=2pt];
		\fill[GeometryBlue] (3.8,3.2) circle[radius=2pt];
		\node[below left] at (0,0) {\((0,0)\)};
		\node[above right] at (3.8,3.2) {\((1,1)\)};
		\node at (1.9,3.85) {\(Q(s,t)=\alpha(s)\beta(t),\quad a=\alpha(1)\)};
		\draw[->,GeometryGray] (-0.8,0)--(-0.8,3.2) node[above] {\(t\)};
		\draw[->,GeometryGray] (0,-0.85)--(3.8,-0.85) node[right] {\(s\)};
	\end{tikzpicture}
	\caption{方形内的对角线与“先沿底边、再沿右边”的折线在固定端点下同伦。图中路径名称表示经 \(Q\) 映射后的路径；因此逐点乘积 \(\alpha\cdot\beta\) 与拼接 \(\alpha*(a\beta)\) 给出同一个道路类。}
	\label{g1:fig:cover-path-product}
\end{figure}


在这个群结构下，基本群的左覆盖变换作用也有直接解释。若 \(k=[\gamma]\in\ker p\)，则对任意 \(q=[\alpha]\)，\eqref{g1:eq:pointwise-product-concatenation}给出
\[
	kq=[\gamma*\alpha].
\]
这恰是第五章前置环路的左作用。因此“左覆盖变换”与“由核元素作左乘”在这里一致；核的中心性又说明它同时等于右乘。

\subsection{离散中心商与同 Lie 代数的不同 Lie 群}
\label{g1:subsec:1604:05}

万有覆盖同态把离散中心核中的元素识别为单位元，却不改变任何点附近的微分结构。反向地，从一个 Lie 群出发作同样的识别，也应得到具有相同 Lie 代数的 Lie 群。此处的离散性让商空间可以由不相交的小邻域直接构造。

\begin{proposition}[离散中心商]
	\label{g1:prop:discrete-central-quotient}
	设 \(G\) 是连通有限维实 Lie 群，\(\Gamma\subset Z(G)\) 为离散子群，其中离散指其子空间拓扑离散。则 \(\Gamma\) 闭，商群 \(G/\Gamma\) 具有唯一的光滑 Lie 群结构，使自然投影 \(q:G\to G/\Gamma\) 为光滑覆盖同态及局部微分同胚。特别地，\(\mathrm{d}q_{e_G}\) 是 Lie 代数同构。
\end{proposition}
\begin{proof}
	先证明闭性。由离散性取单位元开邻域 \(O\)，使 \(O\cap\Gamma=\{e_G\}\)，再取单位元邻域 \(V\)，使 \(V^{-1}V\subset O\)。每个左平移 \(xV\) 至多含一个 \(\Gamma\) 中的点，因为若 \(\gamma_1=xv_1,\gamma_2=xv_2\)，则 \(\gamma_1^{-1}\gamma_2=v_1^{-1}v_2\in O\cap\Gamma\)。若 \(x\in\overline\Gamma\)，开邻域 \(xV\) 必含一个 \(\gamma\in\Gamma\)，并且只含这一个。于是 \(x\) 属于单点 \(\{\gamma\}\) 的闭包；Hausdorff 性使单点闭，所以 \(x=\gamma\in\Gamma\)。故 \(\Gamma\) 闭。

	\(\Gamma\) 位于中心，因而正规，商集上有群乘法。给它商拓扑。投影 \(q\) 开，因为 \(q^{-1}(q(A))=\bigcup_{\gamma\in\Gamma}A\gamma\) 对每个开集 \(A\subset G\) 都开。继续缩小 \(V\)，使它成为坐标邻域且 \(V^{-1}V\cap\Gamma=\{e_G\}\)。则各 \(V\gamma\) 两两不交，且
	\[
		q^{-1}(q(V))=\coprod_{\gamma\in\Gamma}V\gamma.
	\]
	投影在每片上为开连续双射，故为同胚；左平移后，所有 \(q(gV)\) 都有同样的局部形式。因此 \(q\) 是覆盖。

	商空间也是 Hausdorff。若 \(x\Gamma\ne y\Gamma\)，则 \(x^{-1}y\notin\Gamma\)。由 \(\Gamma\) 闭及 \((u,v)\mapsto u^{-1}v\) 连续，可取 \(x,y\) 的开邻域 \(A,B\)，使 \(A^{-1}B\cap\Gamma=\varnothing\)。于是开集 \(q(A),q(B)\) 不交。又因为 \(q\) 开且满，\(G\) 的一个可数基的像给出商空间的可数基。

	把每个 \(q(gV)\) 的坐标从 \(gV\) 传递过去。两套局部逆在交叠点附近相差一个固定的右平移 \(R_\gamma\)：在足够小的连通均匀覆盖邻域上，选定的两张片不能突然改变所对应的 \(\gamma\)。右平移光滑，所以这些坐标相容，并使 \(q\) 局部微分同胚。商乘法的光滑性可在两个因子附近分别选择这样的光滑局部逆，再作 \(G\) 的乘法并投影；商取逆同理。它们与陪集定义一致，因而得到 Lie 群结构。任何使 \(q\) 为局部微分同胚的结构都必须使用同样的局部坐标，故唯一。最后，\(q\) 为光滑同态，微分保括号；局部微分同胚又使其微分可逆，因此得到 Lie 代数同构。
\end{proof}

还需要一个把“微分同构”转换成“覆盖同态”的结论。仅有局部逆函数定理尚未说明映射满射，也未说明所有原像能被同一个小邻域同时控制。

\begin{proposition}[微分同构与覆盖同态]
	\label{g1:prop:lie-homomorphism-local-isomorphism}
	设 \(F:G\to H\) 为连通有限维实 Lie 群之间的光滑同态。若 \(\mathrm{d}F_{e_G}\) 为同构，则 \(F\) 满射且为光滑覆盖映射，\(\ker F\) 是离散中心子群，并诱导 Lie 群同构 \(G/\ker F\cong H\)。
\end{proposition}
\begin{proof}
	逆函数定理给出单位元开邻域 \(U\subset G\)、\(V\subset H\)，使 \(F|_U:U\to V\) 为微分同胚。像 \(F(G)\) 是含 \(V\) 的子群，所以为开子群：对每个像点，左平移 \(V\) 仍留在像中。其余陪集也开，因此像的补集开。由 \(H\) 连通，得到 \(F(G)=H\)。

	记 \(K=\ker F\)。单位元在 \(K\) 中孤立，平移后每个核元素也孤立，所以 \(K\) 离散。它正规；对固定 \(k\in K\)，映射 \(g\mapsto gkg^{-1}\) 连续且取值于离散 \(K\)，源群连通使其恒等于 \(k\)，故 \(K\subset Z(G)\)。

	现在精确写出覆盖片。任意 \(x\in F^{-1}(V)\) 都有唯一 \(u\in U\) 满足 \(F(u)=F(x)\)，于是 \(x=ku\)，其中 \(k=xu^{-1}\in K\)。各 \(kU\) 两两不交：若 \(ku=k'u'\)，则 \(F(u)=F(u')\)，由 \(F|_U\) 单射先得 \(u=u'\)，再得 \(k=k'\)。因此
	\[
		F^{-1}(V)=\coprod_{k\in K}kU,
	\]
	且每片经 \(F\) 微分同胚于 \(V\)。对任意 \(h\in H\)，选 \(g\in G\) 满足 \(F(g)=h\)，左平移这个分解便覆盖 \(hV\)。故 \(F\) 为覆盖。最后，\(F\) 的纤维正是 \(K\) 陪集，因此诱导群双射 \(G/K\to H\)；它在刚构造的小片坐标中是微分同胚，故整体为 Lie 群同构。
\end{proof}

\begin{corollary}[同 Lie 代数的连通 Lie 群]
	\label{g1:cor:connected-groups-same-lie-algebra}
	设 \(G,H\) 是连通有限维实 Lie 群，且 \(\mathfrak g\cong\mathfrak h\)。则存在一个单连通 Lie 群 \(S\) 及其离散中心子群 \(\Gamma_G,\Gamma_H\)，使
	\[
		G\cong S/\Gamma_G,\qquad H\cong S/\Gamma_H.
	\]
	这里可取 \(S\) 为 \(G\) 的万有覆盖群；它同时也是 \(H\) 的一个万有覆盖群。
\end{corollary}
\begin{proof}
	取 \(p:S\to G\) 为\cref{g1:thm:universal-cover-lie-group}构造的万有覆盖同态，并固定 Lie 代数同构 \(\varphi:\mathfrak g\to\mathfrak h\)。复合 \(\varphi\circ\mathrm{d}p_e:\operatorname{Lie}(S)\to\mathfrak h\) 是 Lie 代数同构。因为 \(S\) 单连通，\cref{g1:thm:lie-homomorphism-integration}将它积分为同态 \(F:S\to H\)。由上一命题，\(F\) 是满射覆盖，核 \(\Gamma_H\) 离散且中心。令 \(\Gamma_G=\ker p\)，分别应用上一命题得到两个商同构。\(S\) 单连通，所以 \(F\) 又是 \(H\) 的万有覆盖。
\end{proof}

例如加法群 \(\mathbb R\) 和圆周群 \(S^1\) 都有一维交换 Lie 代数，其共同单连通覆盖群是 \(\mathbb R\)：前者对应核 \(\{0\}\)，后者通过 \(t\mapsto e^{it}\) 对应核 \(2\pi\mathbb Z\)。它们一个非紧、一个紧，因此并不同构。另一方面，\(\mathbb R/(a\mathbb Z)\) 在 \(a>0\) 改变时都同构于圆周；不同的离散子群未必给出不同的群同构类型。

\ProseParagraphBreak

这些结果从已经给定的 Lie 群及其 Lie 代数出发：括号确定局部同态，单连通源群允许唯一的整体积分，离散中心核则记录可以怎样作整体识别。局部微分数据与整体群的关系由此分成了可分别检查的两部分。

% 纲要条目排为编号小节标题，并进入目录与书签；说明段仍保留于纲要注释中。


\begin{chaptersummary}
	单位元切空间经左平移成为左不变向量场，场的括号给出 \(\mathfrak g=\mathrm{T}_eG\) 的李代数结构。共轭微分形成伴随表示 \(\operatorname{Ad}\)，再在单位元求微分便得到 \(\operatorname{ad}_XY=[X,Y]\)。群同态的微分保持括号；矩阵群中，这个括号是 \(XY-YX\)。

	\ProseParagraphBreak

	左不变性让局部积分曲线能够重复平移，得到完备的单参数子群。指数满足 \((\mathrm{d}\exp_G)_0=I\)，故提供单位元附近的坐标；连通群由这样的指数邻域生成，但单个指数未必能覆盖整个群。指数的自然性把同态和共轭连接到李代数，BCH 展开则以 \(\tfrac12[X,Y]\) 为局部乘法偏离加法的首个修正项。

	\ProseParagraphBreak

	李子代数确定左不变的对合分布，过单位元的最大叶形成唯一连通浸入子群。这里必须使用叶自身的流形拓扑。若子群在环境群中闭，指数横截方向不能无限靠近单位元，因而得到局部嵌入结构；无理斜率子群展示了缺少闭性时可能出现的稠密浸入。

	\ProseParagraphBreak

	李代数同态先确定局部群同态。沿路径把小增量依次相乘可延拓其值，同伦网格保证端点固定的同伦不改变结果。源群单连通时得到唯一整体同态。一般连通李群的万有覆盖带有唯一相容李群结构，覆盖核离散且位于中心；同一个已实现李代数的连通李群因而可由其单连通群的不同离散中心商得到。把同态下降到商群时，还必须检查相应离散核是否被保持。
\end{chaptersummary}

\BookExercises

\begin{exercise}\label{g1:ex:affine-group-lie-data}
	在 \(G=(0,\infty)\times\mathbb R\) 上定义
	\[
		(a,b)(c,d)=(ac,b+ad).
	\]
	证明这是一个李群。以单位元处的坐标向量 \(H=\partial_a|_{(1,0)}\)、\(E=\partial_b|_{(1,0)}\) 为基，求它们的左不变延拓与括号，并求 \(\operatorname{Ad}_{(a,b)}(xH+yE)\)。直接对所得伴随表示求导，核对 \(\operatorname{ad}_XY=[X,Y]\)。
\end{exercise}
\begin{solution}
	把 \((a,b)\) 对应到矩阵 \(\begin{pmatrix}a&b\\0&1\end{pmatrix}\)，矩阵相乘恰好给出题设运算，且对应单射，因此运算满足结合律。单位元是 \((1,0)\)，逆为 \((a^{-1},-b/a)\)。\(G\) 是 \(\mathbb R^2\) 的开集，乘法由多项式给出，求逆在 \(a>0\) 时光滑，所以它是二维实李群。这也是实直线上变换 \(t\mapsto at+b\) 的复合群。

	固定 \((a,b)\) 后，左平移为 \(L_{(a,b)}(c,d)=(ac,b+ad)\)。它在单位元的微分把 \((x,y)\) 送到 \((ax,ay)\)，故
	\[
		H^L=a\partial_a,\qquad E^L=a\partial_b.
	\]
	按坐标括号公式，\(H^L\) 对 \(E^L\) 的 \(\partial_b\) 系数 \(a\) 求导得到 \(a\)，而 \(E^L\) 对 \(H^L\) 的 \(\partial_a\) 系数求导为零。因此 \([H^L,E^L]=a\partial_b=E^L\)，即 \([H,E]=E\)。实双线性和反对称性进一步给出
	\[
		[xH+yE,x'H+y'E]=(xy'-yx')E.
	\]

	利用上面的逆元公式，逐项计算共轭得
	\[
		(a,b)(c,d)(a,b)^{-1}=(c,ad+(1-c)b).
	\]
	在 \((c,d)=(1,0)\) 求导，因此
	\[
		\operatorname{Ad}_{(a,b)}(xH+yE)=xH+(ay-bx)E.
	\]
	任取经过单位元的曲线 \((a(t),b(t))\)，若 \(a'(0)=x\)、\(b'(0)=y\)，则对 \(Y=x'H+y'E\)，
	\[
		\left.\frac{\mathrm{d}}{\mathrm{d}t}\right|_{t=0}\operatorname{Ad}_{(a(t),b(t))}Y
		=(xy'-yx')E=[xH+yE,Y].
	\]
	这正是 \(\operatorname{ad}_XY=[X,Y]\)，同时显示了缩放方向与平移方向的括号为何非零：改变缩放会改变平移的大小。
\end{solution}

\begin{exercise}\label{g1:ex:right-invariant-bracket-sign}
	设 \(G\) 为有限维实李群，\(X,Y\in\mathfrak g=\mathrm{T}_eG\)。定义右不变场 \(X^R_g=\mathrm{d}(R_g)_eX\)，左不变延拓仍记为 \(X^L\)。证明求逆微分同胚 \(\iota(g)=g^{-1}\) 满足 \(\iota_*X^L=-X^R\)，并推出
	\[
		[X^R,Y^R]=-[X,Y]^R,\qquad [X^L,Y^R]=0.
	\]
	说明为什么用右不变场定义单位元切空间的括号时，需要调整一个负号才能沿用本书的李代数约定。
\end{exercise}
\begin{solution}
	由乘法光滑，\(g\mapsto\mathrm{d}(R_g)_eX\) 光滑；又由 \(R_a\circ R_g=R_{ga}\)，
	\[
		\mathrm{d}(R_a)_gX^R_g=\mathrm{d}(R_{ga})_eX=X^R_{ga},
	\]
	所以 \(X^R\) 确为右不变场。恒等式 \(\iota\circ L_g=R_{g^{-1}}\circ\iota\) 在 \(e\) 处求导，并用\eqref{g1:eq:lie-group-unit-differentials}的 \(\mathrm{d}\iota_e=-\operatorname{id}\)，得到
	\[
		\mathrm{d}\iota_gX^L_g
		=\mathrm{d}(R_{g^{-1}})_e\mathrm{d}\iota_eX
		=-\mathrm{d}(R_{g^{-1}})_eX=-X^R_{g^{-1}}.
	\]
	这就是 \(\iota_*X^L=-X^R\)。微分同胚保持向量场的括号，故
	\[
		-[X,Y]^R
		=\iota_*([X,Y])^L
		=\iota_*[X^L,Y^L]
		=[-X^R,-Y^R]=[X^R,Y^R].
	\]

	为证明混合括号为零，在单位元取 \(X^L\) 的局部积分曲线 \(\gamma(t)\)。由\eqref{g1:eq:left-invariant-local-flow}，\(X^L\) 的局部流是 \(\Phi_t=R_{\gamma(t)}\)。右不变性意味着 \((R_{\gamma(t)})_*Y^R=Y^R\)，等价地 \(\Phi_t^*Y^R=Y^R\)，在相应的局部定义域内成立。用\cref{g1:prop:flow-pullback-derivative}在 \(t=0\) 求导，得到 \([X^L,Y^R]=0\)。这个论证只需小时间流，不需要先证明左不变场完备。

	最后，第一式在 \(e\) 处给出 \([X^R,Y^R]_e=-[X,Y]\)。因此若希望仍使用由左不变场定义的括号，右不变场版本应写成 \([X,Y]=-[X^R,Y^R]_e\)。负号来自求逆把单位元速度变成相反方向，不能在两种不变性之间直接省去。
\end{solution}

\begin{exercise}\label{g1:ex:so-three-cross-product}
	对 \(v=(v_1,v_2,v_3)\in\mathbb R^3\)，令
	\[
		\widehat v=
		\begin{pmatrix}
			0&-v_3&v_2\\
			v_3&0&-v_1\\
			-v_2&v_1&0
		\end{pmatrix}.
	\]
	证明 \(v\mapsto\widehat v\) 是 \(\mathbb R^3\) 到 \(\mathfrak{so}(3)\) 的线性同构，并满足 \([\widehat v,\widehat w]=\widehat{v\times w}\)。对 \(R\in\operatorname{SO}(3)\)，证明 \(\operatorname{Ad}_R\widehat v=\widehat{Rv}\)。若把 \(R\) 换成任意 \(A\in\operatorname{O}(3)\)，求正确的公式。
\end{exercise}
\begin{solution}
	由\cref{g1:prop:matrix-classical-lie-groups}，\(\mathfrak{so}(3)\) 是全体实反对称 \(3\times3\) 矩阵。任意这样的矩阵恰由三个上三角元素决定，且唯一写成题设形式，所以 \(v\mapsto\widehat v\) 是线性同构。直接乘矩阵可见 \(\widehat vu=v\times u\)。对任意 \(u\in\mathbb R^3\)，展开叉积的三个坐标可得
	\[
		v\times(w\times u)=w(v\cdot u)-u(v\cdot w).
	\]
	交换 \(v,w\) 后相减，便有
	\begin{align*}
		[\widehat v,\widehat w]u
		&=v\times(w\times u)-w\times(v\times u)\\
		&=w(v\cdot u)-v(w\cdot u)=(v\times w)\times u.
	\end{align*}
	因此 \([\widehat v,\widehat w]=\widehat{v\times w}\)。这里左边采用的确实是本书的李代数括号，因为\cref{g1:prop:matrix-lie-bracket}已证明它等于矩阵交换子。

	对正交矩阵 \(A\)，利用叉积的刻画 \((v\times w)\cdot u=\det(v,w,u)\)，可得
	\[
		\bigl((Av)\times(Aw)\bigr)\cdot(Au)
		=\det A\,\det(v,w,u)
		=\bigl(\det A\,A(v\times w)\bigr)\cdot(Au).
	\]
	由于 \(Au\) 遍历所有向量，
	\[
		(Av)\times(Aw)=\det A\,A(v\times w).
	\]
	令 \(w=A^{-1}u\)，并用 \((\det A)^2=1\)，得到
	\[
		(A\widehat v A^{-1})u=A(v\times A^{-1}u)
		=\det A\,(Av)\times u.
	\]
	因而
	\[
		\operatorname{Ad}_A\widehat v=\det A\,\widehat{Av}.
	\]
	在 \(A=R\in\operatorname{SO}(3)\) 时 \(\det R=1\)，便化为题目要求的 \(\operatorname{Ad}_R\widehat v=\widehat{Rv}\)。行列式为 \(-1\) 的正交变换反转空间定向，正是它带来一般公式中的额外负号。
\end{solution}

\begin{exercise}\label{g1:ex:circle-power-homomorphism-kernel}
	设 \(m\) 为非零整数。证明 \(P_m:S^1\to S^1\)、\(P_m(z)=z^m\) 是李群同态，求其单位元微分与核，证明它在每一点都是局部微分同胚，并判断何时为李群同构。再考虑加法李群之间的同态 \(F:\mathbb R\to\mathbb R^2\)、\(F(t)=(t,0)\)：判断 \((0,0)\) 是否为正则值，并直接说明 \(\ker F\) 的光滑结构。这两个例子分别说明了从同态的微分推断整体性质时的什么限制？
\end{exercise}
\begin{solution}
	圆周上的复数乘法交换，故 \((zw)^m=z^mw^m\)，包括负整数 \(m\) 在内。非负整数次幂由复数乘法组成，负整数次幂还要复合光滑求逆，因此 \(P_m\) 光滑，是李群同态。由 \(\mathrm{T}_1S^1=i\mathbb R\)，取经过 \(1\) 的曲线 \(z(t)=\cos t+i\sin t\)，它的初速度为 \(i\)，而 \(P_m(z(t))=\cos(mt)+i\sin(mt)\)，故
	\[
		\mathrm{d}(P_m)_1(it)=imt.
	\]
	\(m\neq0\) 时这个微分可逆，由\cref{g1:prop:lie-homomorphism-constant-rank}，其他点的微分也可逆；逆函数定理说明 \(P_m\) 在每点都是局部微分同胚。

	令 \(q=|m|\)。每个圆周点可写成 \(\cos\theta+i\sin\theta\)，\(z^m=1\) 等价于 \(m\theta\in2\pi\mathbb Z\)。因此
	\[
		\ker P_m=\left\{\cos\frac{2\pi k}{q}+i\sin\frac{2\pi k}{q}:k=0,1,\ldots,q-1\right\}.
	\]
	它是 \(q\) 个点组成的离散李群，其切空间为零，符合 \(\ker\mathrm{d}(P_m)_1=0\)。给定 \(w=\cos\theta+i\sin\theta\)，取 \(z=\cos(\theta/m)+i\sin(\theta/m)\) 即有 \(P_m(z)=w\)，所以 \(P_m\) 总是满射。当 \(|m|>1\) 时核不止一个点，它不是单射；当 \(m=1\) 时是恒等映射，当 \(m=-1\) 时是求逆，二者都有光滑逆。因此恰在 \(m=\pm1\) 时是李群同构。

	对第二个同态，\(\mathrm{d}F_t(v)=(v,0)\) 在每点的秩均为 \(1\)，没有满射到二维靶切空间。唯一原像点为 \(t=0\)，所以 \((0,0)\) 不是正则值。不过 \(\ker F=\{0\}\) 显然是零维嵌入子流形，以唯一的点为零维坐标图即可；其切空间为 \(0=\ker\mathrm{d}F_0\)，也符合恒秩核的结论。

	第一个例子表明微分处处可逆仍不能排除离散的多对一现象，局部微分同胚未必是整体同构。第二个例子表明同态的核是嵌入子流形，并不意味着靶单位元是正则值；在微分不满射时，保证核光滑的是恒秩标准形给出的纤维切片。
\end{solution}

\begin{exercise}\label{g1:ex:affine-exponential}
	令 \(G=\{\left(\begin{smallmatrix}a&b\\0&1\end{smallmatrix}\right):a>0,\ b\in\mathbb R\}\)，采用参数 \((a,b)\) 给出的光滑结构。求其李代数与指数映射，并证明 \(\exp_G:\mathfrak g\to G\) 是整体微分同胚。
\end{exercise}
\begin{solution}
	矩阵乘法和求逆在半平面 \(a>0\) 上光滑，因此 \(G\) 是二维实李群。单位元为 \((1,0)\)，切向矩阵恰为
	\[
		X=\begin{pmatrix}x&y\\0&0\end{pmatrix},\qquad x,y\in\mathbb R.
	\]
	两个这样的矩阵有括号 \([X(x,y),X(u,v)]=X(0,xv-uy)\)，从而给出这二维空间的完整李代数结构。对 \(k\geq1\)，\(X^k=\left(\begin{smallmatrix}x^k&x^{k-1}y\\0&0\end{smallmatrix}\right)\)。由矩阵指数级数及包含同态的自然性，
	\[
		\exp_GX=\begin{pmatrix}e^x&y\psi(x)\\0&1\end{pmatrix},\qquad
		\psi(x)=\int_0^1e^{tx}\,\mathrm{d}t
		=\begin{cases}(e^x-1)/x,&x\ne0,\\1,&x=0.\end{cases}
	\]
	积分表达式说明 \(\psi\) 在整个实轴光滑且严格为正，因而分段写法在零点没有奇性。给定 \((a,b)\) 且 \(a>0\)，唯一原像为
	\[
		x=\log a,\qquad y=\frac{b}{\psi(\log a)}.
	\]
	这也光滑，所以指数是整体微分同胚。这里全局可逆是明确计算的结果，不能从任意李群的单位元附近局部可逆性直接推出。
\end{solution}

\begin{exercise}\label{g1:ex:non-surjective-exponential}
	在 \(\operatorname{SL}(2,\mathbb R)\) 中令 \(A=\operatorname{diag}(-2,-1/2)\)。证明 \(A\) 属于单位连通分支，且是两个指数的乘积，但不存在实矩阵 \(X\) 满足 \(e^X=A\)。这说明连通李群的指数不一定满射。
\end{exercise}
\begin{solution}
	令
	\[
		J=\begin{pmatrix}0&-1\\1&0\end{pmatrix},\qquad D=\begin{pmatrix}1&0\\0&-1\end{pmatrix}.
	\]
	二者迹为零，属于 \(\mathfrak{sl}(2,\mathbb R)\)。直接求指数得 \(e^{\pi J}=-I\)、\(e^{(\log2)D}=\operatorname{diag}(2,1/2)\)，故
	\[
		A=e^{\pi J}e^{(\log2)D}.
	\]
	曲线 \(t\mapsto e^{t\pi J}e^{t(\log2)D}\)，\(0\leq t\leq1\)，留在 \(\operatorname{SL}(2,\mathbb R)\) 中，并从 \(I\) 走到 \(A\)。所以 \(A\) 属于单位连通分支 \(G^0\)，两个因子的整条单参数子群也都在 \(G^0\) 中。

	若实矩阵 \(X\) 满足 \(e^X=A\)，则 \(X\) 与 \(A\) 交换，因为 \(X\) 与自身的每个幂交换，进而与收敛的指数级数交换。写出 \(XA=AX\) 的两个非对角元，由 \(-2\ne-1/2\) 得 \(X_{12}=X_{21}=0\)。于是 \(X\) 为实对角矩阵，其指数的两个对角元都严格为正，不可能等于 \(A\)。特别地，连通李群 \(G^0\) 的指数不能到达 \(A\)，但有限个指数的乘积仍能到达它。
\end{solution}

\begin{exercise}\label{g1:ex:commutator-exponential-limit}
	设 \(G\) 为有限维实李群，\(X,Y\in\mathfrak g\)。对充分小的实数 \(s,t\)，令
	\[
		K(s,t)=\exp_G(sX)\exp_G(tY)\exp_G(-sX)\exp_G(-tY).
	\]
	证明 \(\log K(s,t)=st[X,Y]+O((|s|+|t|)^3)\)，并由此证明
	\[
		K(1/n,1/n)^{n^2}\longrightarrow\exp_G([X,Y]).
	\]
	说明改变交换子的次序为何会改变首项的符号。
\end{exercise}
\begin{solution}
	由共轭与指数的相容性，前面三个因子的乘积为
	\[
		\exp_G\bigl(t e^{s\operatorname{ad}_X}Y\bigr).
	\]
	固定 \(X,Y\)，记 \(\delta=|s|+|t|\)，并令 \(U=t e^{s\operatorname{ad}_X}Y\)、\(V=-tY\)。算子指数展开给出
	\[
		U=tY+st[X,Y]+O(\delta^3),\qquad U+V=st[X,Y]+O(\delta^3).
	\]
	\([tY,-tY]=0\)，故 \([U,V]=O(\delta^3)\)。将它们代入 BCH 的二阶展开，得到所需对数公式；此处 \(K\) 对充分小的参数确在单位元对数邻域内。

	令 \(Z_n=\log K(1/n,1/n)\)。则 \(Z_n=n^{-2}[X,Y]+O(n^{-3})\)，所以 \(n^2Z_n\to[X,Y]\)。单参数子群规律给出 \(K(1/n,1/n)^{n^2}=\exp_G(n^2Z_n)\)，连续性即得极限。若把 \(X,Y\) 两方向的角色交换，首项变为 \(st[Y,X]=-st[X,Y]\)；这里的符号由写出的乘积顺序决定。
\end{solution}

\begin{exercise}\label{g1:ex:adjoint-kernel-center}
	设 \(G\) 为连通有限维实李群，\(Z(G)=\{g\in G:gh=hg\text{ 对所有 }h\in G\}\)。证明 \(\ker\operatorname{Ad}=Z(G)\)。再证明若 \([X,Y]=0\) 对所有 \(Y\in\mathfrak g\) 成立，则 \(\exp_GX\in Z(G)\)，并给出逆命题不成立的例子。
\end{exercise}
\begin{solution}
	若 \(g\in Z(G)\)，共轭 \(C_g\) 是恒等映射，故其单位元微分为恒等，即 \(g\in\ker\operatorname{Ad}\)。反之，若 \(\operatorname{Ad}_g=I\)，指数自然性说明
	\[
		g\exp_G(Y)g^{-1}=\exp_G(Y)\qquad(Y\in\mathfrak g).
	\]
	因 \(G\) 连通，每个群元素都是有限个指数的乘积，故 \(g\) 与所有群元素交换。这一步正是连通性的用途。

	若 \(\operatorname{ad}_X=0\)，则 \(\operatorname{Ad}_{\exp_GX}=e^{\operatorname{ad}_X}=I\)，于是 \(\exp_GX\) 为中心元素。反向可能因指数的周期而失败。取 \(G=\operatorname{SL}(2,\mathbb R)^0\)，并令
	\[
		X=\pi\begin{pmatrix}0&-1\\1&0\end{pmatrix},\qquad
		Y=\begin{pmatrix}1&0\\0&-1\end{pmatrix}.
	\]
	\(X,Y\in\operatorname{Lie}(G)=\mathfrak{sl}(2,\mathbb R)\)，其中等号来自单位分支为开子群。\(\exp_GX=-I\) 与所有矩阵交换，但
	\[
		[X,Y]=2\pi\begin{pmatrix}0&1\\1&0\end{pmatrix}\ne0.
	\]
	因此“某一个指数落入中心”不能替代“整个方向与所有方向的括号为零”。
\end{solution}

\begin{exercise}\label{g1:ex:closed-subgroups-real-line}
	证明加法李群 \((\mathbb R,+)\) 的闭子群恰为 \(\{0\}\)、\(\mathbb R\) 以及 \(a\mathbb Z\)、\(a>0\)。求它们的李代数与单位元连通分支，并说明为什么李代数不能区分所有闭子群。
\end{exercise}
\begin{solution}
	设 \(K\subseteq\mathbb R\) 为闭子群。若 \(K\) 没有正元素，由取负封闭可知它也没有负元素，故 \(K=\{0\}\)。以下设正元素集非空，并令
	\[
		a=\inf\{x\in K:x>0\}\ge0.
	\]

	若 \(a>0\)，由下确界的定义可取 \(x_j\in K\)、\(x_j>0\)，使 \(x_j\to a\)。闭性给出 \(a\in K\)。任取 \(x\in K\)，选择整数 \(m\) 使 \(ma\le x<(m+1)a\)，则 \(r=x-ma\in K\) 且 \(0\le r<a\)。若 \(r>0\)，就与 \(a\) 是正元素下确界矛盾，所以 \(r=0\)，即 \(x\in a\mathbb Z\)。反向包含由 \(a\in K\) 和子群性质得到，故 \(K=a\mathbb Z\)。

	若 \(a=0\)，可取 \(0<x_j\in K\) 且 \(x_j\to0\)。对任意 \(y\in\mathbb R\)，选整数 \(m_j\) 使
	\[
		|m_jx_j-y|\le x_j.
	\]
	于是 \(m_jx_j\in K\) 且趋于 \(y\)，闭性迫使 \(y\in K\)。因此 \(K=\mathbb R\)。反过来，\(\{0\}\) 和 \(\mathbb R\) 显然闭；当 \(a>0\) 时，\(a\mathbb Z\) 的每个点与其余点距离至少为 \(a\)，而不在该集合中的点位于两个相邻整数倍之间，存在避开它的开区间，所以 \(a\mathbb Z\) 也闭。

	加法群的指数为恒等映射，故闭子群定理给出的李代数为
	\[
		\mathfrak k=\{X\in\mathbb R:tX\in K\text{ 对所有 }t\in\mathbb R\}.
	\]
	对 \(K=\mathbb R\)，它等于 \(\mathbb R\)；对 \(K=\{0\}\) 或 \(K=a\mathbb Z\)，它等于 \(\{0\}\)，因为非零 \(X\) 的全部实数倍充满 \(\mathbb R\)。相应的单位元连通分支分别是 \(\mathbb R\)、\(\{0\}\) 和 \(\{0\}\)。所以所有这些非平凡离散子群都具有同一个零李代数；不同的离散分支不会出现在单位元切空间中。
\end{solution}

\begin{exercise}\label{g1:ex:triangular-two-dimensional-subgroups}
	考虑上三角矩阵李群
	\[
		N=\left\{
		n(x,y,z)=
		\begin{pmatrix}
			1&x&z\\
			0&1&y\\
			0&0&1
		\end{pmatrix}:x,y,z\in\mathbb R
		\right\}.
	\]
	记其李代数的基为 \(X=E_{12}\)、\(Y=E_{23}\)、\(Z=E_{13}\)。证明每个二维李子代数都恰为
	\[
		\mathfrak h_{a,b}=\operatorname{span}\{aX+bY,Z\},
		\qquad(a,b)\ne(0,0).
	\]
	求出相应的连通积分子群，并判断它是否闭。说明为什么 \(\operatorname{span}\{X,Y\}\) 不能成为二维浸入李子群的李代数。
\end{exercise}
\begin{solution}
	矩阵乘法与交换子分别给出
	\[
		n(x,y,z)n(x',y',z')
		=n(x+x',y+y',z+z'+xy'),
	\]
	以及
	\[
		[X,Y]=Z,\qquad[X,Z]=[Y,Z]=0.
	\]
	所以对 \(U=aX+bY+cZ\)、\(V=a'X+b'Y+c'Z\)，有
	\[
		[U,V]=(ab'-ba')Z.
	\]

	设 \(\mathfrak h\) 为二维李子代数。若 \(Z\notin\mathfrak h\)，则投影
	\[
		\pi:\mathfrak h\longrightarrow\operatorname{span}\{X,Y\},
		\qquad aX+bY+cZ\longmapsto aX+bY
	\]
	的核为零，所以它在两个二维空间之间是同构。取 \(\mathfrak h\) 的一组基 \(U,V\)，其投影独立，因而 \(ab'-ba'\ne0\)。括号封闭性随即给出 \(Z\in\mathfrak h\)，矛盾。因此 \(Z\in\mathfrak h\)。再选一个不与 \(Z\) 共线的元素 \(aX+bY+cZ\)，减去 \(cZ\) 即得所述形式，其中 \((a,b)\ne(0,0)\)。反过来，任意这样的平面内两个元素的括号都为零，故确为李子代数。

	令
	\[
		H_{a,b}=\{n(x,y,z):bx-ay=0\}
		=\{n(as,bs,z):s,z\in\mathbb R\}.
	\]
	乘法公式说明它对乘积封闭。逆矩阵为
	\[
		n(x,y,z)^{-1}=n(-x,-y,xy-z),
	\]
	所以它也对取逆封闭。约束 \(bx-ay=0\) 是闭的线性条件，故 \(H_{a,b}\) 在 \(N\) 中闭。参数 \((s,z)\mapsto n(as,bs,z)\) 是从 \(\mathbb R^2\) 到这个平面的微分同胚，说明它连通；在单位元求导得到切空间 \(\operatorname{span}\{aX+bY,Z\}\)。由\cref{g1:thm:lie-subalgebra-integration}的唯一性，它就是所求的连通积分子群。这里的参数 \(z\) 是矩阵右上角的实际条目，无须先把它换成指数坐标。

	最后，\([X,Y]=Z\notin\operatorname{span}\{X,Y\}\)，所以这个平面不是李子代数。浸入李子群的单位切空间必须对括号封闭，故不可能由任何二维浸入李子群实现。两个切向方向逐次运动产生的第三方向，正是本例的障碍。
\end{solution}

\begin{exercise}\label{g1:ex:vector-line-stabilizer-subgroups}
	设 \(n\ge2\)，\(v\in\mathbb R^n\setminus\{0\}\)。在 \(\operatorname{GL}(n,\mathbb R)\) 中定义
	\[
		K_v=\{g:gv=v\},\qquad
		K_{\mathbb Rv}=\{g:g(\mathbb Rv)=\mathbb Rv\}.
	\]
	证明这两个集合都是闭子群，分别求其李代数和维数。这里的闭性相对于 \(\operatorname{GL}(n,\mathbb R)\)，不要求它们在所有 \(n\times n\) 矩阵组成的向量空间中闭。
\end{exercise}
\begin{solution}
	保持一个向量或一条直线的可逆线性映射，其复合与逆仍保持相应对象，所以两集合都是子群。选一组以 \(v\) 为首向量的基。在这组基中，它们分别具有块矩阵形式
	\[
		K_v=
		\left\{
		\begin{pmatrix}
			1&a\\
			0&B
		\end{pmatrix}:
		\begin{aligned}
			&a\in\operatorname{Mat}_{1,n-1}(\mathbb R),\\
			&B\in\operatorname{GL}(n-1,\mathbb R)
		\end{aligned}
		\right\},
	\]
	\[
		K_{\mathbb Rv}=
		\left\{
		\begin{pmatrix}
			c&a\\
			0&B
		\end{pmatrix}:
		\begin{aligned}
			&c\in\mathbb R\setminus\{0\},\\
			&a\in\operatorname{Mat}_{1,n-1}(\mathbb R),\\
			&B\in\operatorname{GL}(n-1,\mathbb R)
		\end{aligned}
		\right\}.
	\]
	前者由第一列等于 \((1,0,\ldots,0)^{\mathsf T}\) 的闭条件描述；后者由第一列的后 \(n-1\) 项为零的闭条件描述。后一情形中，\(c\ne0\) 与 \(B\) 可逆已由整个矩阵可逆保证。因此两者在 \(\operatorname{GL}(n,\mathbb R)\) 中闭，可以使用\cref{g1:thm:closed-subgroup}。

	若 \(\exp(tA)v=v\) 对所有 \(t\) 成立，在 \(t=0\) 求导得到 \(Av=0\)。反之，若 \(Av=0\)，则 \(A^kv=0\) 对所有 \(k\ge1\) 成立，指数级数给出 \(\exp(tA)v=v\)。于是
	\[
		\operatorname{Lie}(K_v)=\{A:Av=0\}
		=
		\left\{
		\begin{pmatrix}
			0&a\\
			0&C
		\end{pmatrix}
		\right\}.
	\]

	若 \(\exp(tA)v\in\mathbb Rv\) 对所有 \(t\) 成立，对商空间投影 \(\mathbb R^n\to\mathbb R^n/\mathbb Rv\) 求导可得 \(Av\in\mathbb Rv\)。反之，若 \(Av=\lambda v\)，则指数级数给出 \(\exp(tA)v=e^{t\lambda}v\)，它始终张成原来的直线。故
	\[
		\operatorname{Lie}(K_{\mathbb Rv})=\{A:Av\in\mathbb Rv\}
		=
		\left\{
		\begin{pmatrix}
			\lambda&a\\
			0&C
		\end{pmatrix}
		\right\}.
	\]
	这里 \(C\) 是任意 \((n-1)\times(n-1)\) 实矩阵。计数自由条目，得到
	\[
		\dim K_v=(n-1)+(n-1)^2=n(n-1),\qquad
		\dim K_{\mathbb Rv}=n^2-n+1.
	\]
	相差的一个维数来自沿指定直线的伸缩：固定向量禁止这一变化，保持直线则允许它。
\end{solution}

\begin{exercise}\label{g1:ex:connected-subgroup-normal-in-closure}
	设 \(H\) 为有限维实李群 \(G\) 的连通浸入李子群，\(\mathfrak h=\operatorname{Lie}(H)\)，令 \(K=\overline H\)，并赋予它闭子群定理给出的李群结构，记 \(\mathfrak k=\operatorname{Lie}(K)\)。证明
	\[
		\operatorname{Ad}_k\mathfrak h=\mathfrak h\quad(k\in K),
		\qquad[\mathfrak k,\mathfrak h]\subseteq\mathfrak h,
	\]
	且 \(H\) 是 \(K\) 的正规子群。第二个包含式称 \(\mathfrak h\) 为 \(\mathfrak k\) 的理想。用无理斜率环面子群说明，结论并不意味着 \(\mathfrak h=\mathfrak k\)。
\end{exercise}
\begin{solution}
	包含映射 \(H\to G\) 的像落在嵌入子流形 \(K\) 中，由\cref{g1:prop:embedded-initial-property}，它作为到 \(K\) 的映射仍光滑。因而单位元微分给出 \(\mathfrak h\subseteq\mathfrak k\)。对 \(h\in H\)，内部共轭是 \(H\) 的光滑自同构，包含到 \(G\) 后与环境共轭相容，所以
	\[
		\operatorname{Ad}_h\mathfrak h=\mathfrak h.
	\]

	任取 \(k\in K\)。流形在每点有可数邻域基，由闭包的定义可以选 \(h_j\in H\) 使 \(h_j\to k\)，这里的收敛在环境群中理解。对 \(Y\in\mathfrak h\)，伴随作用连续，所以
	\[
		\operatorname{Ad}_{h_j}Y\longrightarrow\operatorname{Ad}_kY.
	\]
	有限维向量空间的线性子空间是闭集，而左侧每一项均属于 \(\mathfrak h\)，故 \(\operatorname{Ad}_kY\in\mathfrak h\)。再对 \(k^{-1}\) 作同一论证，得到反向包含，证明 \(\operatorname{Ad}_k\mathfrak h=\mathfrak h\)。

	若 \(X\in\mathfrak k\)，则闭子群定理保证 \(\exp_G(tX)\in K\) 对所有实数 \(t\) 成立。因此对 \(Y\in\mathfrak h\)，上一结论给出
	\[
		\operatorname{Ad}_{\exp_G(tX)}Y\in\mathfrak h.
	\]
	在零点取差商，使用\eqref{g1:eq:ad-bracket}以及 \(\mathfrak h\) 的闭线性性质，得到 \([X,Y]\in\mathfrak h\)。这就是理想条件。

	最后固定 \(k\in K\)。共轭 \(C_k\) 是 \(G\) 的微分同胚，把 \(H\) 的流形结构传到 \(kHk^{-1}\)，便得到一个连通浸入李子群，其李代数为 \(\operatorname{Ad}_k\mathfrak h=\mathfrak h\)。由\cref{g1:thm:lie-subalgebra-integration}，同一李子代数的连通积分子群唯一，所以 \(kHk^{-1}=H\)。这对每个 \(k\in K\) 成立，即 \(H\) 在 \(K\) 中正规；这里使用的是 \(H\) 自身的连通性。

	在\cref{g1:exm:irrational-torus-subgroup}中，\(H=\{[t,\alpha t]:t\in\mathbb R\}\)，\(\alpha\) 无理，且 \(K=T^2\)。于是 \(\mathfrak h=\mathbb R(1,\alpha)\) 真包含于 \(\mathfrak k=\mathbb R^2\)。环面群交换，两者的括号都为零，任何线性子空间都是理想。这同时满足本题所有结论，并显示闭包确实可以增大李代数。
\end{solution}

\begin{exercise}\label{g1:ex:torus-homomorphism-integer-matrix}
	记 \(T^m=\mathbb R^m/\mathbb Z^m\)，用商映射的单位元微分把其李代数识别为 \(\mathbb R^m\)，\(T^n\) 同理。设 \(A:\mathbb R^m\to\mathbb R^n\) 为实线性映射。证明 \(A\) 能积分为李群同态 \(T^m\to T^n\)，当且仅当 \(A\) 在标准基下的全部矩阵条目都是整数；存在时写出这个同态。对 \(m=n=1\)，说明非零整数 \(k\) 所给同态的每个纤维有多少点，并说明 \(k=0\) 的情形。
\end{exercise}
\begin{solution}
	设 \(q_m:\mathbb R^m\to T^m\)、\(q_n:\mathbb R^n\to T^n\) 为商同态。若 \(F:T^m\to T^n\) 的单位元微分为 \(A\)，则 \(F\circ q_m\) 与 \(q_n\circ A\) 都是从连通群 \(\mathbb R^m\) 到 \(T^n\) 的同态，且微分相同。由\cref{g1:thm:lie-homomorphism-integration}的唯一性，
	\[
		F(q_m(x))=q_n(Ax).
	\]
	对每个 \(z\in\mathbb Z^m\)，左边在 \(x=z\) 时为单位元，所以 \(Az\in\mathbb Z^n\)。特别地，对标准基向量逐一使用这个条件，得到 \(A\) 的每一列都是整数向量，即所有矩阵条目为整数。

	反过来，若 \(A\mathbb Z^m\subset\mathbb Z^n\)，定义 \(F([x])=[Ax]\)。改变代表 \(x\) 为 \(x+z\) 时，目标改变 \(Az\in\mathbb Z^n\)，所以映射良定义。它保持加法，并在商覆盖的局部坐标中由线性映射 \(A\) 表示，因而光滑，微分也正是 \(A\)。

	一维时映射为 \([x]\mapsto[kx]\)。若 \(k\ne0\)，其核由 \(kx\in\mathbb Z\) 决定，共有 \(|k|\) 个模 \(\mathbb Z\) 的不同解。映射满射，因为对任意代表 \(y\in\mathbb R\) 可取 \(x=y/k\)；每个纤维是核的一个陪集，故也有 \(|k|\) 个点。\(k=0\) 时则为常值同态，单位元的纤维是整个圆周，其他点的纤维为空，因而不是到 \(T^1\) 的覆盖。
\end{solution}

\begin{exercise}\label{g1:ex:connected-abelian-lie-group-quotient}
	设 \(G\) 是 \(n\) 维连通实李群，且其李代数括号恒为零。利用同态积分证明 \(G\) 为交换群，并证明存在加法群 \(\mathbb R^n\) 的离散子群 \(\Gamma\)，使 \(G\cong\mathbb R^n/\Gamma\)。这里已经给定 \(G\)，不要求使用交换李群的分类定理。
\end{exercise}
\begin{solution}
	加法群 \(\mathbb R^n\) 的李代数也是 \(\mathbb R^n\)，括号恒为零。选线性同构 \(L:\mathbb R^n\to\mathfrak g\)，因为两边括号都为零，它是李代数同构。源群 \(\mathbb R^n\) 连通且单连通，所以\cref{g1:thm:lie-homomorphism-integration}将 \(L\) 积分成同态 \(F:\mathbb R^n\to G\)。

	\(\mathrm{d}F_0=L\) 为同构，目标 \(G\) 连通，因此\cref{g1:prop:lie-homomorphism-local-isomorphism}使 \(F\) 成为满射覆盖，核 \(\Gamma\) 离散，并给出 \(G\cong\mathbb R^n/\Gamma\)。加法群 \(\mathbb R^n\) 交换，其任意同态像都交换；由满射性，\(G\) 交换。也可在商的表示中直接读出这一点。\(n=0\) 时，连通零维群只有一个点，同一结论取 \(\Gamma=\{0\}\)。
\end{solution}

\begin{exercise}\label{g1:ex:discrete-normal-subgroup-central}
	证明连通李群的任意离散正规子群都位于中心。再在 \(\operatorname{O}(2)\) 中取旋转 \(r=R_{2\pi/3}\)，证明 \(\Gamma=\{I,r,r^2\}\) 是离散正规子群，但不位于中心，由此说明连通条件的作用。
\end{exercise}
\begin{solution}
	设 \(G\) 连通、\(\Gamma\) 离散且正规。固定 \(\gamma\in\Gamma\)。正规性保证连续映射
	\[
		G\longrightarrow\Gamma,\qquad g\longmapsto g\gamma g^{-1}
	\]
	确实取值于 \(\Gamma\)。连通空间的连续像连通，而离散空间中连通子集至多含一点，所以该映射恒等于它在单位元的值 \(\gamma\)。因此 \(g\gamma=\gamma g\) 对所有 \(g\) 成立，即 \(\Gamma\subset Z(G)\)。

	现在令
	\[
		R_\theta=\begin{pmatrix}\cos\theta&-\sin\theta\\
		\sin\theta&\cos\theta\end{pmatrix},\qquad
		S=\begin{pmatrix}1&0\\0&-1\end{pmatrix}.
	\]
	每个二维正交矩阵都是某个 \(R_\theta\) 或 \(R_\theta S\)：第一列为单位向量 \((\cos\theta,\sin\theta)\)，第二列因正交性及单位长度只能是 \(\pm(-\sin\theta,\cos\theta)\)。旋转彼此交换，直接相乘又得 \(SR_\theta S^{-1}=R_{-\theta}\)。所以任意正交矩阵共轭 \(\Gamma\) 时，都将它保持不变，\(\Gamma\) 正规。它有限，故在 Hausdorff 矩阵空间中离散。

	但 \(SrS^{-1}=r^{-1}=r^2\ne r\)，所以 \(r\) 不在中心。共轭映射在 \(\operatorname{O}(2)\) 的旋转分支与反射分支上可以取不同值；源空间不连通，便不能用前段的常值论证把两者连在一起。
\end{solution}

\begin{exercise}\label{g1:ex:lie-homomorphism-descent}
	设 \(S\) 为单连通实李群，\(\Gamma\subset Z(S)\) 为离散子群，\(q:S\to S/\Gamma\) 为商同态。给定有限维实李群 \(H\) 及李代数同态 \(\varphi:\operatorname{Lie}(S)\to\mathfrak h\)，令 \(\widetilde F:S\to H\) 为其整体积分。证明存在同态 \(F:S/\Gamma\to H\) 满足 \(F\circ q=\widetilde F\)，当且仅当 \(\Gamma\subset\ker\widetilde F\)，并说明存在时为何光滑且唯一。将结论用于 \(S=(\mathbb R,+)\)、\(\Gamma=2\pi\mathbb Z\)、\(H=(\mathbb R,+)\) 及 \(\varphi(t)=\lambda t\)，判断哪些 \(\lambda\) 可以下降。
\end{exercise}
\begin{solution}
	若已有 \(F\)，则对 \(\gamma\in\Gamma\)，有
	\[
		\widetilde F(\gamma)=F(q(\gamma))=F(e_{S/\Gamma})=e_H,
	\]
	所以 \(\Gamma\subset\ker\widetilde F\)。

	反过来，若这个包含成立，定义 \(F(s\Gamma)=\widetilde F(s)\)。当 \(s'=s\gamma\) 时，
	\(\widetilde F(s')=\widetilde F(s)\widetilde F(\gamma)=\widetilde F(s)\)，故定义无关于陪集代表；同态性质由 \(\widetilde F\) 传下。由\cref{g1:prop:discrete-central-quotient}，\(q\) 在每个点附近有光滑局部逆 \(\sigma\)。在相应商邻域上，
	\[
		F=\widetilde F\circ\sigma,
	\]
	所以 \(F\) 光滑。投影 \(q\) 满射，等式 \(Fq=\widetilde F\) 又唯一决定 \(F\)。

	在最后给定的加法群例子中，\(\widetilde F(t)=\lambda t\)。下降要求 \(\widetilde F(2\pi)=2\pi\lambda=0\)，所以只有 \(\lambda=0\) 可以下降。于是任意光滑同态 \(S^1\cong\mathbb R/(2\pi\mathbb Z)\to(\mathbb R,+)\) 都是零同态。与圆周作为目标时的整数条件相比，这里目标没有可容纳非零一周增量的周期识别。
\end{solution}

\begin{exercise}\label{g1:ex:local-isomorphism-simply-connected-target}
	设 \(F:G\to H\) 是连通有限维实李群之间的光滑同态，且 \(\mathrm{d}F_{e_G}\) 为同构。证明：若目标群 \(H\) 单连通，则 \(F\) 为李群同构。再给出目标群不是单连通时结论失败的例子，并核对其单位元微分。
\end{exercise}
\begin{solution}
	令 \(\varphi=(\mathrm{d}F_{e_G})^{-1}:\mathfrak h\to\mathfrak g\)。它是李代数同态，因为李代数同构的逆保持括号。目标群 \(H\) 现在作为源群，具有单连通性，所以\cref{g1:thm:lie-homomorphism-integration}给出 \(K:H\to G\)，满足 \(\mathrm{d}K_{e_H}=\varphi\)。

	链式法则说明 \(K\circ F:G\to G\) 的单位元微分为恒等，而恒等同态的微分也是恒等。源群 \(G\) 连通，同态唯一性给出 \(KF=\operatorname{id}_G\)。同样，\(FK=\operatorname{id}_H\)。因此 \(F\) 有光滑同态逆 \(K\)，为李群同构。这个证明没有预先假定 \(G\) 单连通；它作为最终与 \(H\) 微分同胚的结果才随之成立。

	反例为
	\[
		F:(\mathbb R,+)\longrightarrow S^1,\qquad F(t)=e^{it}.
	\]
	单位元处的切向量 \(1\in \mathrm{T}_0\mathbb R\) 被送成 \(i\in \mathrm{T}_1S^1\)，按单位角速度识别两边的李代数后，微分就是 \(\mathbb R\) 上的恒等映射。但 \(\ker F=2\pi\mathbb Z\)，故 \(F\) 不单射。它是覆盖同态；非平凡的离散核正是局部同构没有变成整体同构的原因。
\end{solution}
